\SourcePage{579}
\SGAsection{4}{特征 $p>0$ 所特有的现象}
在整个第 4 节中，$p$ 表示一个素数，$k$ 表示\Term{characteristic}{特征}为 $p$ 的\Term{pseudocompact-ring}{伪紧环}，$\pi$ 表示 $k$ 的\Term{continuous}{连续}\Term{endomorphism}{自同态} $x\mapsto x^p$。

\SourcePage{580}
\textbf{4.1。} 设 $X$ 是\Term{affine-algebra}{仿射代数}为 $A$ 的 $k$-\Term{formal-variety}{形式簇}。把经\Term{base-change}{基变换} $\pi:k\to k$ 得到的 $k$-\Term{formal-variety}{形式簇} $X\widehat\otimes_\pi k$（见 1.2.D）记为 $X^{(p/k)}$ 或 $X^{(p)}$；其\Term{affine-algebra}{仿射代数}是\Term{completed-tensor-product}{完备张量积} $A\widehat\otimes_\pi k$。\Term{continuous}{连续}\Term{morphism}{态射} $A\widehat\otimes_\pi k\to A$，$a\widehat\otimes_\pi\lambda\mapsto a^p\lambda$，诱导一个 $k$-\Term{formal-variety}{形式簇}\Term{morphism}{态射}
\[
\operatorname{Fr}(X/k):X\longrightarrow X^{(p/k)}.
\]
以下称 $\operatorname{Fr}(X/k)$ 为 $X$ 相对于 $k$ 的 Frobenius \Term{morphism}{态射}，并常简记为 $\operatorname{Fr}$。

\textbf{4.1.1。}\footnote{编者注（134）：这里已修正原文；原文把 $\widehat{X^{(p)}}/\widehat S\subset(\widehat X/\widehat S)^{(p)}$ 的包含关系写反了。另外，本段在后文中并未用到。}
现在考虑\Term{prime-field}{素域} $\mathbf F_p$ 上的\Term{scheme}{概形} $S$ 及一个 $S$-\Term{scheme}{概形} $\eta:X\to S$。设 $\operatorname{fr}(S)$ 为 $S$ 的“绝对”Frobenius \Term{morphism}{态射}（它在\Term{underlying-topological-space}{底拓扑空间}上诱导\Term{identity-map}{恒等映射}，并把\Term{structure-sheaf}{结构层}的任一\Term{section}{截面} $f$ 映为 $f^p$；见第 VII$_A$ 讲 4.1），$X^{(p)}$ 为 $S$-\Term{scheme}{概形} $X\times_{\operatorname{fr}(S)}S$（第 VII$_A$ 讲，同上），即\Term{structure-morphism}{结构态射} $X^{(p)}\to S$ 是 $\operatorname{fr}(S)\circ\eta$。

设 $\widehat S$ 是这样的\Term{formal-scheme}{形式概形}：其\Term{underlying-topological-space}{底拓扑空间}是 $S$ 的\Term{closed-point}{闭点}集，赋予\Term{discrete-topology}{离散拓扑}；在\Term{closed-point}{闭点} $s$ 处的\Term{local-ring}{局部环} $\mathcal O_{\widehat S,s}$ 是 $\mathcal O_{S,s}$ 关于\Term{finite}{有限}余\Term{length}{长度}\Term{ideal}{理想}定义的\Term{linear-topology}{线性拓扑}的\Term{separated-completion}{分离完备化} $\widehat{\mathcal O}_{S,s}$（见 1.2.6）。因此，其\Term{affine-algebra}{仿射代数} $k=\mathcal A(\widehat S)$ 是这些 $\widehat{\mathcal O}_{S,s}$ 的乘积，$s$ 遍历 $S$ 的\Term{closed-point}{闭点}。回顾同一处的记号：$\widehat X/\widehat S$ 表示由满足 $[\kappa(x):\kappa(s)]<\infty$ 的点 $x\in X$ 组成的 $k$-\Term{formal-variety}{形式簇}，其中 $s$ 是 $x$ 在 $S$ 中的像；其\Term{local-ring}{局部环} $\mathcal O_{\widehat X/\widehat S,x}$ 是 $\mathcal O_{X,x}$ 关于下述\Term{ideal}{理想} $I$ 定义的\Term{linear-topology}{线性拓扑}的\Term{separated-completion}{分离完备化} $\widehat{\mathcal O}_{X,x}$：$\mathcal O_{X,x}/I$ 作为 $\mathcal O_{X,x}$-模具有\Term{finite-length}{有限长度}（因而作为 $\mathcal O_{S,s}$-模也有\Term{finite-length}{有限长度}）。由此可经\Term{base-change}{基变换}构造 $k$-\Term{formal-variety}{形式簇}
\[
(\widehat X/\widehat S)^{(p)}=(\widehat X/\widehat S)\widehat\otimes_\pi k.
\]
也可考虑 $k$-\Term{formal-variety}{形式簇} $\widehat{X^{(p)}}/\widehat S$：其\Term{underlying-set}{底集}由这样的 $x\in X^{(p)}=X$ 组成，记 $s$ 为 $x$ 在 $S$ 中的像，则\Term{morphism}{态射}
\[
\begin{tikzcd}\kappa(s)\arrow[r,"\operatorname{fr}"]&\kappa(s)\arrow[r,"\eta^\natural"]&\kappa(x)\end{tikzcd}
\]
使 $\kappa(x)$ 成为 $\kappa(s)$ 的\Term{finite}{有限}次\Term{group-extension}{扩张}。这时 $\eta^\natural:\kappa(s)\to\kappa(x)$ 也具有这一性质，即 $x$ 是 $\widehat X/\widehat S$ 的点，而且
\[
\mathcal O_{\widehat{X^{(p)}}/\widehat S,x}
=\widehat{\mathcal O}_{X,x}\widehat\otimes_\pi\widehat{\mathcal O}_{S,s}
=\widehat{\mathcal O}_{X^{(p)},x}.
\]
第二个等号成立，是因为 $X\mapsto\widehat X/\widehat S$ 与乘积可交换，见 1.2.6。因此，$\widehat{X^{(p)}}/\widehat S$ 可等同于 $(\widehat X/\widehat S)^{(p)}$ 的一个\Term{formal-subvariety}{形式子簇}。并且，二者相等当且仅当对 $S$ 的每个\Term{closed-point}{闭点} $s$，$\kappa(s)$ 在 $\kappa(s)^p$ 上为\Term{finite}{有限}次\Term{group-extension}{扩张}。例如，若 $S$ 在一个满足 $[\kappa:\kappa^p]<+\infty$ 的域 $\kappa$ 上\Term{locally-of-finite-type}{局部有限型}，则此条件成立。

\SourcePage{581}
\textbf{4.1.2。} 设 $G$ 为一个 $k$-\Term{formal-group}{形式群}。\Term{functor}{函子} $X\mapsto X^{(p)}=X\widehat\otimes_\pi k$ 与\Term{finite}{有限}乘积可交换，因而把 $k$-\Term{formal-group}{形式群}变为 $k$-\Term{formal-group}{形式群}。而且 $\operatorname{Fr}(X/k)$ 关于 $X$ 是\Term{functor}{函子}性的，所以\Term{morphism}{态射}
\[\operatorname{Fr}:G\longrightarrow G^{(p)}\]
是 $k$-\Term{formal-group}{形式群}\Term{homomorphism}{同态}。若置 $G^{(p^{n+1})}=(G^{(p^n)})^{(p)}$，则复合\Term{morphism}{态射}
\[
\operatorname{Fr}^n=\operatorname{Fr}^n(G/k):
G\xrightarrow{\operatorname{Fr}}G^{(p)}\xrightarrow{\operatorname{Fr}}G^{(p^2)}\xrightarrow{\operatorname{Fr}}\cdots\xrightarrow{\operatorname{Fr}}G^{(p^n)}
\]
也如此。记 $A$ 为 $G$ 的\Term{affine-algebra}{仿射代数}，$I_A$ 为其\Term{augmentation-ideal}{增广理想}。

\begin{SGAresult}{定义}{4.1.2}
(a) 把 $\operatorname{Fr}^n$ 的核记为 ${}_{\operatorname{Fr}^n}G$。直接由定义可知，${}_{\operatorname{Fr}^n}G$ 是\Term{infinitesimal}{无穷小}的，其\Term{affine-algebra}{仿射代数}为商 $A/I_A^{\{p^n\}}$，其中 $I_A^{\{p^n\}}$ 表示由 $I_A$ 中元素的 $p^n$ 次幂生成的 $A$ 的\Term{closed-ideal}{闭理想}。

(b) 若 $I_A^{\{p^n\}}=0$，即 ${}_{\operatorname{Fr}^n}G=G$，则称 $G$ 的\Term{height}{高度} $\leq n$。
\end{SGAresult}

\textbf{4.2。} 显然，$k$-\Term{formal-group}{形式群} $G$ 的 \Term{lie-algebra}{Lie 代数} $\operatorname{Lie}(G)$ 是\Term{algebra}{代数} $\mathrm H(G)$ 的一个 \Term{p-lie-subalgebra}{$p$-Lie 子代数}（见 2.3）。事实上，立即可归结到 $k$ 为 \Term{artinian-ring}{阿廷环}的情形；此时沿用 2.3 和 2.6 (c) 的记号，$\operatorname{Lie}(G)$ 是 $\mathrm H(G)$ 中满足 $\varphi_G(x\otimes1+1\otimes x)=\Delta_G(x)$ 的元素 $x$ 所成的部分。而
\[
\begin{aligned}
\varphi_G(x^p\otimes1+1\otimes x^p)
&=\varphi_G\bigl((x\otimes1+1\otimes x)^p\bigr)\\
&=\varphi_G(x\otimes1+1\otimes x)^p
=\Delta_G(x)^p=\Delta_G(x^p).
\end{aligned}
\]

\textbf{4.2.1。} 反过来，$\mathcal O_k$ 上的每个 \Term{p-lie-algebra}{$p$-Lie 代数} $L$ 都定义一个 $k$-\Term{group-functor}{群值函子}。事实上，记 $\mathrm U_p(L)$ 为把 $\mathbf{Alf}/k$ 的\Term{object}{对象} $C$ 映为 $C$ 上的 \Term{p-lie-algebra}{$p$-Lie 代数} $L(C)$ 的\Term{restricted-enveloping-algebra}{限制包络代数} $\mathrm U_p(L(C))$ 的\Term{functor}{函子}（第 VII$_A$ 讲 5.3）。由第 VII$_A$ 讲 5.4，$\mathrm U_p(L)$ 是一个 $\mathcal O_k$-\Term{bialgebra}{双代数}，故依 2.2 确定一个 $k$-\Term{group-functor}{群值函子} $\operatorname{Spf}^{*}\mathrm U_p(L)$，今后记为 $\mathscr G_p(L)$。

因此，$\mathscr G_p(L)(C)$ 是 $\mathrm U_p(L(C))$ 中\Term{augmentation}{增广}为 $1$ 且满足 $\Delta_{\mathrm U_p(L(C))}(z)=z\otimes z$ 的元素 $z$ 所成的群。

\textbf{4.2.2。} 设 $L$ 是 $\mathcal O_k$ 上的\Term{flat}{平坦} \Term{p-lie-algebra}{$p$-Lie 代数}。结合第 VII$_A$ 讲 5.3.3，可像命题 2.6.2 (i) 那样证明 $\mathrm U_p(L)$ 在 $\mathcal O_k$ 上\Term{flat}{平坦}；因此 $\mathscr G_p(L)$ 是一个\Term{topologically-flat}{拓扑平坦}的 $k$-\Term{formal-group}{形式群}，其\Term{covariant}{协变} $\mathcal O_k$-\Term{bialgebra}{双代数}是 $\mathrm U_p(L)$。

由 2.6.2 (iii) 的证明，对每个\Term{finite-length}{有限长度} $k$-\Term{algebra}{代数} $C$，$\operatorname{Lie}(\mathscr G_p(L))(C)$ 是 $\mathrm U_p(L)(C)=\mathrm U_p(L(C))$ 的\Term{primitive-element}{本原元}所成的\Term{set}{集合}（另见第 VII$_A$ 讲 3.2.3）；再由第 VII$_A$ 讲 5.5.3，\Term{canonical}{典范}\Term{morphism}{态射} $L\to\mathrm U_p(L)$ 诱导 \Term{p-lie-algebra}{$p$-Lie 代数}\Term{isomorphism}{同构}\footnote{编者注（135）：这里详述了以下论证。}
\[
\tau_{p,L}:L\xrightarrow{\sim}\operatorname{Lie}(\mathscr G_p(L))
\]
（与\Term{characteristic}{特征} $0$ 时的 3.1 比较）。

\SourcePage{582}
\Term{formal-group}{形式群} $\mathscr G_p(L)$ 可由一个\Term{universal-property}{泛性质}刻画。事实上，$\mathscr G_p(L)$ 到\Term{formal-group}{形式群} $G$ 的每个\Term{morphism}{态射} $h$ 都诱导一个\Term{morphism}{态射} $\operatorname{Lie}(h):\operatorname{Lie}(\mathscr G_p(L))\to\operatorname{Lie}(G)$。将它与\Term{isomorphism}{同构} $\tau_{p,L}$ 复合，得到\Term{morphism}{态射} $h':L\to\operatorname{Lie}(G)$。

\begin{SGAresult}{命题}{4.2.2}
若 $k$ 是\Term{characteristic}{特征} $p>0$ 的\Term{pseudocompact-ring}{伪紧环}，$G$ 是一个 $k$-\Term{formal-group}{形式群}，$L$ 是 $\mathcal O_k$ 上的\Term{flat}{平坦} \Term{p-lie-algebra}{$p$-Lie 代数}，则上述映射 $h\mapsto h'$ 是一个\Term{bijection}{双射}
\[
\operatorname{Hom}_{\mathbf{Grf}/k}(\mathscr G_p(L),G)
\xrightarrow{\sim}\operatorname{Hom}_{p\text{-}\mathrm{Lie}}(L,\operatorname{Lie}(G)),
\]
右端表示 $L$ 到 $\operatorname{Lie}(G)$ 的 \Term{p-lie-algebra}{$p$-Lie 代数}\Term{morphism}{态射}集。
\end{SGAresult}
\begin{proof}
立即可归结到 $k$ 为 \Term{artinian-ring}{阿廷环}的情形。\footnote{编者注（136）：证明与 2.6.3 相同。}置 $\mathfrak L=L(k)$。由 2.3.1，$\operatorname{Hom}_{\mathbf{Grf}/k}(\mathscr G_p(L),G)$ 与使下图交换的含幺\Term{algebra}{代数}\Term{morphism}{态射} $h:\mathrm U_p(\mathfrak L)\to\mathrm H(G)$ 所成的\Term{set}{集合}\Term{bijection}{双射}：
\[
\begin{tikzcd}[column sep=large]
\mathrm U_p(\mathfrak L)\arrow[r,"h"]\arrow[d,"\Delta_{\mathrm U_p(\mathfrak L)}"']&\mathrm H(G)\arrow[dr,"\delta_G"]&\\
\mathrm U_p(\mathfrak L)\otimes\mathrm U_p(\mathfrak L)\arrow[r,"h\otimes h"']&\mathrm H(G)\otimes\mathrm H(G)\arrow[r,"\varphi_G"']&\mathrm H(G\times G).
\end{tikzcd}
\]
而 $h$ 由其在 $\mathfrak L$ 上的\Term{restriction}{限制}确定；该\Term{restriction}{限制}是 $\mathfrak L$ 到 $\mathrm H(G)$ 的底 \Term{p-lie-algebra}{$p$-Lie 代数}的\Term{morphism}{态射}。图的\Term{commutativity}{交换性}意味着 $h$ 把 $\mathfrak L$ 映入 $\mathrm H(G)$ 中满足 $\delta_G(x)=\varphi_G(x\otimes1+1\otimes x)$ 的元素所成的部分，而这一部分正是 $\operatorname{Lie}(G)$，见 2.6 (c)。
\end{proof}

\textbf{4.3。} 现在我们更详细地研究 $L$ 为 $\mathcal O_k$ 上\Term{flat}{平坦} \Term{p-lie-algebra}{$p$-Lie 代数}时的 $k$-\Term{formal-group}{形式群} $\mathscr G_p(L)$。

为此，先考虑\Term{characteristic}{特征} $p$ 的环 $C$ 及 $C$ 上一个 \Term{p-lie-algebra}{$p$-Lie 代数} $\mathfrak L$，其底模是以 $(x_i)_{i\in I}$ 为基的\Term{free-module}{自由模}。又记 $P$ 为满足 $0\leq n_i<p$ 且仅有\Term{finite}{有限}多个非零项的非负整数族 $\mathbf n=(n_i)_{i\in I}$ 的\Term{set}{集合}。在 $I$ 上选取一个\Term{totally-ordered}{全序}，若 $i_1<i_2<\cdots<i_r$ 是使 $n_i\ne0$ 的指标，则置
\[
|\mathbf n|=n_{i_1}+\cdots+n_{i_r},\qquad
x^{\mathbf n}=x_{i_1}^{n_{i_1}}x_{i_2}^{n_{i_2}}\cdots x_{i_r}^{n_{i_r}},\qquad
\mathbf n!=(n_{i_1}!)\cdots(n_{i_r}!).
\]
已知单项式 $x^{\mathbf n}/\mathbf n!$ 构成 $\mathrm U_p(\mathfrak L)$ 的一组基（第 VII$_A$ 讲 5.3.3）；显然有
\[
\Delta\left(\frac{x^{\mathbf n}}{\mathbf n!}\right)
=\sum_{\mathbf m\leq\mathbf n}\frac{x^{\mathbf m}}{\mathbf m!}\otimes
\frac{x^{\mathbf n-\mathbf m}}{(\mathbf n-\mathbf m)!}.\tag{*}
\]
\SourcePage{583}
其中求和遍历满足 $\mathbf m\leq\mathbf n$（即对一切 $i\in I$，$m_i\leq n_i$）的所有 $\mathbf m\in P$。

由公式 (*) 容易确定 $\mathrm U_p(\mathfrak L)$ 的自然\Term{filtration}{滤过}（见 1.3.6）。置 $\mathrm U=\mathrm U_p(\mathfrak L)$，设 $\mathrm U_+$ 为由 $\mathfrak L$ 生成的\Term{two-sided-ideal}{双边理想}，$\mathrm U_0=C\cdot1_{\mathrm U}$。如 1.3.6，对 $n\geq1$ 置
\[
\mathrm U_n=\{u\in\mathrm U\mid\Delta_{\mathrm U}(u)-u\otimes1\in\mathrm U_{n-1}\otimes\mathrm U_+\},
\]
得到 $\mathrm U$ 的一个由 $C$-\Term{subcoalgebra}{子余代数}组成的\Term{filtration}{滤过}。公式 (*) 表明 $\mathrm U_n$ 是以满足 $|\mathbf m|\leq n$ 的单项式 $x^{\mathbf m}$ 为基的自由 $C$-模。

\textbf{4.3.1。} 沿用 4.3 的记号，确定 $\mathrm U=\mathrm U_p(\mathfrak L)$ 中满足 $\varepsilon_{\mathrm U}(\xi)=1$ 和 $\Delta_{\mathrm U}(\xi)=\xi\otimes\xi$ 的元素 $\xi$。任一 $\xi\in\mathrm U$ 可写成 $\xi=\sum_{\mathbf n\in P}\xi_{\mathbf n}(x^{\mathbf n}/\mathbf n!)$，其中 $\xi_{\mathbf n}\in C$。条件 $\varepsilon_{\mathrm U}(\xi)=1$ 意味着 $\xi_{\mathbf0}=1$，$\mathbf0$ 表示所有\Term{set-component}{分量}均为零的 $P$ 的元素。条件 $\Delta_{\mathrm U_p(\mathfrak L)}(\xi)=\xi\otimes\xi$ 意味着
\[
\sum_{\mathbf m\leq\mathbf n}\xi_{\mathbf n}\frac{x^{\mathbf m}}{\mathbf m!}\otimes
\frac{x^{\mathbf n-\mathbf m}}{(\mathbf n-\mathbf m)!}
=\sum_{\mathbf q,\mathbf r}\xi_{\mathbf q}\xi_{\mathbf r}
\frac{x^{\mathbf q}}{\mathbf q!}\otimes\frac{x^{\mathbf r}}{\mathbf r!},
\]
即当所有 $q_i+r_i<p$ 时 $\xi_{\mathbf q+\mathbf r}=\xi_{\mathbf q}\xi_{\mathbf r}$，否则 $\xi_{\mathbf q}\xi_{\mathbf r}=0$。若把 $n_i=1$ 且 $j\ne i$ 时 $n_j=0$ 的 $\xi_{\mathbf n}$ 记为 $\xi_i$，则这些条件\Term{equivalence}{等价}于
\[
\xi_{\mathbf n}=\prod_i\xi_i^{n_i}\quad(\mathbf n\in P),\qquad \xi_i^p=0\quad(i\in I).
\]
由此得到：
\begin{SGAresult}{命题}{4.3.1}
设 $k$ 为\Term{characteristic}{特征} $p>0$ 的\Term{local-scheme}{局部}\Term{pseudocompact-ring}{伪紧环}，\footnote{编者注（137）：这里增加了 $k$ 为\Term{local-ring}{局部环}的假设，以保证每个\Term{topologically-flat}{拓扑平坦}的\Term{pseudocompact}{伪紧} $k$-模均\Term{topologically-free}{拓扑自由}，见证明。}$L$ 为 $\mathcal O_k$ 上的\Term{flat}{平坦} \Term{p-lie-algebra}{$p$-Lie 代数}，$C$ 为 $\mathbf{Alf}/k$ 的\Term{object}{对象}，$\sqrt[p]{0_C}$ 为 $C$ 中满足 $x^p=0$ 的元素构成的\Term{ideal}{理想}。则存在一个“关于 $C$ \Term{functor}{函子}性”的\Term{bijection}{双射}
\[L(C)\otimes_C\sqrt[p]{0_C}\xrightarrow{\sim}\mathscr G_p(L)(C).\]
\end{SGAresult}
\begin{proof}
由注记 1.2.3.F，可为 $L(C)$ 选取一组 $C$-基 $({}_Cx_i)_{i\in I}$，使得对 $\mathbf{Alf}/k$ 的任一\Term{morphism}{态射} $\varphi:C\to D$，$L(\varphi)$ 把 ${}_Cx_i$ 映为 ${}_Dx_i$。根据以上论证，映射
\[
\sum_i{}_Cx_i\otimes\xi_i\longmapsto
\sum_{\mathbf n\in P}\left(\prod_i\xi_i^{n_i}\right)\frac{{}_Cx^{\mathbf n}}{\mathbf n!}
\]
便是所需的\Term{functor}{函子}性\Term{bijection}{双射}。
\end{proof}

\SourcePage{584}
\textbf{4.3.2。} “在\Term{characteristic}{特征} $p>0$ 时不存在 Campbell--Hausdorff 公式。”具体说来：4.3.1 的\Term{functorial-isomorphism}{函子性同构}依赖于基 $({}_Cx_i)_{i\in I}$ 的选取。与 3.2（\Term{characteristic}{特征} $0$ 的情形）不同，一般不存在“同时关于 $C$ 和 $L$ \Term{functor}{函子}性”的\Term{bijection}{双射} $L(C)\otimes_C\sqrt[p]{0_C}\to\mathscr G_p(L)(C)$。以下论证例如表明，若 $k$ 是一个包含所有 $(p^2-1)$ 次\Term{root-of-unity}{单位根}的域，则这样的\Term{isomorphism}{同构}不存在。

选取 $L$，使得对每个 $C\in\mathbf{Alf}/k$，$L(C)$ 为由一个满足关系 $x^{(p^2)}=0$ 的元素 $x$ 生成的交换 \Term{p-lie-algebra}{$p$-Lie 代数}。于是
\[
L(C)=Cx\oplus Cx^{(p)},\qquad\mathrm U_p(L(C))\simeq k[x]/(x^{p^2}).
\]
我们将证明：与 $L$ 的\Term{endomorphism}{自同态}相容且把 $L(C)\otimes_C\sqrt[p]{0_C}$ 映入 $\mathscr G_p(L)(C)$ 的唯一\Term{functor}{函子}性\Term{morphism}{态射}
\[\chi_C:L(C)\otimes_C\sqrt[p]{0_C}\longrightarrow\mathrm U_p(L(C))\]
是取常值 $1$ 的\Term{morphism}{态射}。

令 $\psi_C$ 为 4.3.1 给出的 $L(C)\otimes_C\sqrt[p]{0_C}$ 到 $\mathscr G_p(L)(C)=\operatorname{Prim}\mathrm U_p(L(C))$ 的\Term{bijection}{双射}，即
\[
x\otimes a+x^{(p)}\otimes b\longmapsto\sum_{0\leq i,j<p}a^ib^jx^{i+pj}.
\]
将 $\chi_C$ 与 $\psi_C^{-1}$ 复合，得到关于 $C$ 的\Term{functor}{函子}性\Term{morphism}{态射}
\[
\begin{aligned}
\theta_C:L(C)\otimes_C\sqrt[p]{0_C}&\longrightarrow L(C)\otimes_C\sqrt[p]{0_C},\\
x\otimes a+x^{(p)}\otimes b&\longmapsto x\otimes P(a,b)+x^{(p)}\otimes Q(a,b).
\end{aligned}
\]
关于 $C$ 的\Term{functor}{函子}性表明，$P(a,b)$ 和 $Q(a,b)$ 是两个\Term{polynomial}{多项式} $P,Q\in k[x,y]$ 在 $(a,b)$ 的取值，可设它们关于 $X$ 和 $Y$ 的\Term{degree}{次数}均小于 $p$。\footnote{编者注（138）：以下详述了原文的论证。}对 $\alpha\in k$，令 $\ell_\alpha$ 为把 $x$ 映为 $\alpha x$（因而把 $x^{(p)}$ 映为 $\alpha^px^{(p)}$）的 $L$ 的 \Term{p-lie-algebra}{$p$-Lie 代数}\Term{endomorphism}{自同态}。于是 $\mathrm U_p(\ell_\alpha)$ 是把 $x$ 映为 $\alpha x$（因而对 $i<p^2$，把 $x^i$ 映为 $\alpha^ix^i$）的\Term{algebra}{代数}\Term{endomorphism}{自同态}。容易看出下方块交换：
\[
\begin{tikzcd}
L(C)\otimes_C\sqrt[p]{0_C}\arrow[r,"\psi_C"]\arrow[d,"\ell_\alpha(C)\otimes_C\mathrm{id}"']&\mathrm U_p(L(C))\arrow[d,"\mathrm U_p(\ell_\alpha)(C)"]\\
L(C)\otimes_C\sqrt[p]{0_C}\arrow[r,"\psi_C"']&\mathrm U_p(L(C)).
\end{tikzcd}
\]
假设 $\chi_C\circ(\ell_\alpha(C)\otimes\mathrm{id})=\mathrm U_p(\ell_\alpha)(C)\circ\chi_C$ 因而\Term{pushout}{推出}
\[
P(\alpha a,\alpha^pb)=\alpha P(a,b),\qquad
Q(\alpha a,\alpha^pb)=\alpha^pQ(a,b).
\]
写成 $P=\sum_{i,j<p}\lambda_{ij}X^iY^j$ 及 $Q=\sum_{i,j<p}\mu_{ij}X^iY^j$，并取 $C=k[X,Y]/(X^p,Y^p)$，可知若 $\lambda_{ij}\ne0$（相应地，$\mu_{ij}\ne0$），则 $\alpha^{i-1+pj}=1$（相应地，$\alpha^{i+p(j-1)}=1$）。
\SourcePage{585}
取 $\alpha$ 为一个 $p^2-1$ 阶\Term{primitive-root-of-unity}{本原单位根}，便知 $P=\lambda X$、$Q=\mu Y$，其中 $\lambda,\mu\in k$。因此
\[
\chi_C(x\otimes a+x^{(p)}\otimes b)=\sum_{0\leq i,j<p}(\lambda a)^i(\mu b)^jx^{i+pj}.
\]
最后考虑 $L$ 中把 $x$ 映为 $x+x^{(p)}$ 的\Term{endomorphism}{自同态} $f$；取 $C=k[a]/(a^3)$，比较 $(\chi_C\circ f(C))(x\otimes a)$ 与 $(\mathrm U_p(f)(C)\circ\chi_C)(x\otimes a)$ 中 $x^p$ 及 $x^{p+1}$ 的系数，得到 $\lambda=\mu$ 及 $\lambda^2a^2=0$，故 $\lambda=\mu=0$。\footnote{编者注（138，续）：这里详述了原文的论证。}

\begin{SGAresult}{\Term{theorem}{定理}}{4.4}
设 $k$ 为\Term{characteristic}{特征} $p>0$ 的\Term{pseudocompact-ring}{伪紧环}，$G$ 为一个 $k$-\Term{formal-group}{形式群}。以下断言\Term{equivalence}{等价}：

(i) 存在 $\mathcal O_k$ 上的\Term{flat}{平坦} \Term{p-lie-algebra}{$p$-Lie 代数} $L$，使 $G$ \Term{isomorphism}{同构}于 $\mathscr G_p(L)$（此时由 4.2.2，$L=\operatorname{Lie}(G)$）。

(ii) 存在一个\Term{projective-object}{投射}\Term{pseudocompact}{伪紧} $k$-模 $\omega$，使 $G$ 的\Term{affine-algebra}{仿射代数}\Term{isomorphism}{同构}于 $k[[\omega]]$ 模掉由 $x^p$（$x\in\omega$）生成的\Term{closed-ideal}{闭理想}所得的商（此时 $\omega\simeq\omega_{G/k}$）。

(iii) $G$ 的\Term{height}{高度} $\leq1$，且 $\omega_{G/k}$ 是一个\Term{projective-object}{投射}\Term{pseudocompact}{伪紧} $k$-模。

(iv) $G$ 的\Term{height}{高度} $\leq1$，且在 $k$ 上\Term{topologically-flat}{拓扑平坦}。
\end{SGAresult}
\begin{proof}
(i)$\Rightarrow$(ii)：置 $\omega=\Gamma^*(L)$（见 1.2.3.D），令 $A$ 为 $k[[\omega]]$ 模掉由 $x^p$（$x\in\omega$）生成的\Term{closed-ideal}{闭理想}所得的商。对 $\mathbf{Alf}/k$ 的\Term{object}{对象} $C$，每个\Term{continuous}{连续}\Term{morphism}{态射} $h:A\to C$ 由其在 $\omega$ 上的\Term{restriction}{限制} $h'$ 确定；$h'$ 把 $\omega$ 映入 $\sqrt[p]{0_C}$。由此得到 $\operatorname{Hom}_{\mathbf{Alp}/k}(A,C)$ 到\Term{continuous}{连续} $k$-\Term{linear}{线性}映射集 $\operatorname{Hom}_c(\omega,\sqrt[p]{0_C})$ 的\Term{canonical}{典范}\Term{bijection}{双射}。后者\Term{canonical}{典范}\Term{isomorphism}{同构}于 $L(C)\otimes_C\sqrt[p]{0_C}$（见 3.3 的证明）。于是 4.3.1 所得的\Term{functor}{函子}性\Term{bijection}{双射}
\[L(C)\otimes_C\sqrt[p]{0_C}\xrightarrow{\sim}\mathscr G_p(L)(C)\]
给出 (i)$\Rightarrow$(ii)。

其余蕴含关系可参看\Term{theorem}{定理} 3.3 及第 VII$_A$ 讲 7.4.1 的类似证明。
\end{proof}

\textbf{注记 4.4.A。}\footnote{编者注（139）：这条注记为编者所加，将用于 4.4.2。}设 $G$ 为一个\Term{affine-algebra}{仿射代数}是 $A$ 的\Term{infinitesimal}{无穷小} $k$-\Term{formal-group}{形式群}，并且 $\omega_{G/k}=I_A/\overline{I_A^2}$ 是\Term{projective-object}{投射}\Term{pseudocompact}{伪紧} $k$-模。则\Term{projection}{投影} $I_A\to\omega_{G/k}$ 有一个\Term{section}{截面} $\tau:\omega_{G/k}\to I_A$，而 $\tau$ 诱导一个满的\Term{continuous}{连续}\Term{algebra}{代数}\Term{morphism}{态射} $\psi:k[[\omega_{G/k}]]\to A$，见 3.3 中 (iii)$\Rightarrow$(i) 的证明。

\begin{SGAresult}{推论}{4.4.1}
若 $k$ 为\Term{characteristic}{特征} $p>0$ 的域，则\Term{functor}{函子} $L\mapsto\mathscr G_p(L)$ 给出 $k$ 上 \Term{p-lie-algebra}{$p$-Lie 代数}\Term{category}{范畴}到\Term{height}{高度} $\leq1$ 的 $k$-\Term{formal-group}{形式群}\Term{category}{范畴}的一个\Term{equivalence}{等价}。
\end{SGAresult}
\SourcePage{586}
\begin{proof}
事实上，由\Term{theorem}{定理} 4.4，当 $G$ 遍历\Term{height}{高度} $\leq1$ 的\Term{formal-group}{形式群}时，\Term{functor}{函子} $G\mapsto\mathscr G_p(\operatorname{Lie}G)$ \Term{isomorphism}{同构}于\Term{identity-functor}{恒等函子}；同样，4.2.2 已表明\Term{functor}{函子} $L\mapsto\operatorname{Lie}\mathscr G_p(L)$ \Term{isomorphism}{同构}于\Term{identity-functor}{恒等函子}（另见第 VII$_A$ 讲 5.5.3）。\footnote{编者注（140）：若 $k$ 是\Term{characteristic}{特征} $p>0$ 的 \Term{artinian-ring}{阿廷环}，同一证明给出 $k$ 上\Term{flat}{平坦} \Term{p-lie-algebra}{$p$-Lie 代数}\Term{category}{范畴}与在 $k$ 上\Term{topologically-flat}{拓扑平坦}、\Term{height}{高度} $\leq1$ 的 $k$-\Term{formal-group}{形式群}\Term{category}{范畴}之间的\Term{equivalence}{等价}。}
\end{proof}

\textbf{4.4.2。} 仍设 $k$ 为\Term{characteristic}{特征} $p$ 的域。令 $G$ 为\Term{affine-algebra}{仿射代数}是 $A$ 的\Term{infinitesimal}{无穷小} $k$-\Term{formal-group}{形式群}。因 $G$ \Term{infinitesimal}{无穷小}，$A$ 的每个开\Term{ideal}{理想}在 $n$ \Term{large-relative-to-cardinal}{足够大}时都包含 $I_A^{\{p^n\}}$，所以 $G$ 是群 ${}_{\operatorname{Fr}^n}G$ 的\Term{affine-algebra}{仿射代数} $A/I_A^{\{p^n\}}$ 的\Term{projective-limit}{投射极限}（见 4.1.2）。因此，每个\Term{infinitesimal}{无穷小} $k$-\Term{formal-group}{形式群}都是\Term{finite}{有限}\Term{height}{高度}的 $k$-\Term{formal-group}{形式群}的\Term{inductive-limit}{归纳极限}。

设 $G$ 的\Term{height}{高度} $\leq n$，记 $H=G/{}_{\operatorname{Fr}}G$。\footnote{编者注（141）：以下详述了原文的论证。}由 2.4 和 2.4.1，$\operatorname{Fr}:G\to G^{(p)}$ \Term{resolution}{分解}为\Term{epimorphism}{满态射} $\pi:G\to H$ 后接\Term{monomorphism}{单态射} $i:H\to G^{(p)}$。于是有\Term{commutative-diagram}{交换图}
\[
\begin{tikzcd}
G\arrow[r,"\pi"]\arrow[dr,"\operatorname{Fr}"']&H\arrow[r,"\operatorname{Fr}^{n-1}"]\arrow[d,"i"]&H^{(p^{n-1})}\arrow[d,"i^{(p^{n-1})}"]\\
&G^{(p)}\arrow[r,"\operatorname{Fr}^{n-1}"']&G^{(p^n)}.
\end{tikzcd}
\]
因\Term{functor}{函子} $X\mapsto X^{(p)}$ 与\Term{fiber-product}{纤维积}可交换，$i^{(p^{n-1})}$ 仍为\Term{monomorphism}{单态射}。因 $\operatorname{Fr}^n(G/k)$ 为零且 $\pi$ 为\Term{epimorphism}{满态射}，$\operatorname{Fr}^{n-1}(G^{(p)}/k)\circ i=i^{(p^{n-1})}\circ\operatorname{Fr}^{n-1}(H/k)$ 为零；再由 $i^{(p^{n-1})}$ 为\Term{monomorphism}{单态射}，$\operatorname{Fr}^{n-1}(H/k)$ 为零，即 $H$ 的\Term{height}{高度} $\leq n-1$。因此，每个\Term{finite}{有限}\Term{height}{高度}的 $k$-\Term{formal-group}{形式群}都有一个\Term{composition-series}{合成列}，其商的\Term{height}{高度}均 $\leq1$，因而可用 $k$ 上的 \Term{p-lie-algebra}{$p$-Lie 代数}描述。

最后，$G$ 的\Term{affine-algebra}{仿射代数} $A$ 是 $k[[\omega_{G/k}]]$ 的商，见 4.4.A。因此，若 $\omega_{G/k}$ 在 $k$ 上\Term{finite-dimensional}{有限维}，则所有\Term{algebra}{代数} $A/I_A^{\{p^n\}}$ 都是\Term{finite-dimensional}{有限维} $k$-\Term{vector-space}{向量空间}。故对\Term{characteristic}{特征} $p>0$ 的域上的每个\Term{infinitesimal}{无穷小}\Term{formal-group}{形式群} $G$，若 $\omega_{G/k}$ 在 $k$ 上\Term{finite-dimensional}{有限维}，则 $G$ 是\Term{finite}{有限}\Term{formal-group}{形式群}（即\Term{finite-length}{有限长度}的\Term{formal-group}{形式群}，见 1.2.6）的\Term{inductive-limit}{归纳极限}。

\SGAsection{5}{域上的无穷小形式群的齐性空间}
\textbf{5.0。}\footnote{编者注（142）：本小节为编者所加，以\Term{cocommutative}{余交换} \Term{hopf-algebra}{Hopf 代数}的语言表述后面的命题 5.1，并介绍此后在这一方向上所得的结果。}为简便起见，设 $k$ 为域。令 $G$ 为\Term{affine-algebra}{仿射代数}是 $A=\mathcal A(G)$ 的 $k$-\Term{formal-group}{形式群}。设 $A^+$ 为 $A$ 的\Term{augmentation-ideal}{增广理想}；对 $a\in A$，以 $\overline a=a-\varepsilon_A(a)1_A$ 表示它在 $A^+$ 上的\Term{projection}{投影}。
\SourcePage{587}
设 $F$ 为由 $A$ 的\Term{closed-ideal}{闭理想} $J$ 定义的 $G$ 的\Term{formal-subgroup}{形式子群}，$\pi$ 为\Term{projection}{投影} $A\to A/J$，$\overline\pi$ 为\Term{projection}{投影} $A\to A^+\to A^+/J$ 的复合。注意，对每个 $a\in A$，$\Delta_A(a)$ 在 $A\widehat\otimes A^+$ 上的\Term{projection}{投影}是 $\Delta_A(a)-a\widehat\otimes1$。

由\Term{theorem}{定理} 2.4，\Term{constructible}{可构造}商 $k$-\Term{formal-variety}{形式簇} $X=G/F$，其\Term{affine-algebra}{仿射代数} $B$ 为
\[
\begin{aligned}
B&=\operatorname{Ker}\left(A\mathrel{\substack{\xrightarrow{\tau_1}\\[-2pt]\xrightarrow[(\mathrm{id}_A\widehat\otimes\pi)\Delta_A]{}}}A\widehat\otimes_k(A/J)\right)\\
&=\{a\in A\mid(\mathrm{id}_A\widehat\otimes\pi)\Delta_A(a)=a\widehat\otimes\pi(1)\}\\
&=\operatorname{Ker}\bigl((\mathrm{id}_A\widehat\otimes\overline\pi)\Delta_A\bigr).
\end{aligned}
\]
它也就是由满足以下条件的元素 $\phi$ 组成的 $A$ 的\Term{subalgebra}{子代数}：对所有 $C\in\operatorname{Ob}(\mathbf{Alf}/k)$、$g\in G(C)$、$h\in F(C)$，都有 $\phi(gh)=\phi(g)$。对任意 $g,g'\in G(C)$ 和 $h\in F(C)$，有 $\phi(gg'h)=\phi(gg')$，故 $\Delta(\phi)$ 属于 $\mathrm{id}_A\widehat\otimes_k(\mathrm{id}_A\widehat\otimes\overline\pi)\Delta_A$ 的核；因 $A$ 在 $k$ 上\Term{topologically-flat}{拓扑平坦}，该核等于 $A\widehat\otimes_k B$。于是 $\Delta_A(B)\subset A\widehat\otimes_k B$，即闭\Term{subalgebra}{子代数} $B$ 同时也是一个\Term{left-coideal}{左余理想}。

另一方面，$B$ 确定 $F$，因为由推论 2.4.1，$J=\overline{AB^+}$，即 $J$ 是由 $B^+=B\cap A^+$ 生成的\Term{closed-ideal}{闭理想}。因此得到一个\Term{injectivity}{单射} $Q:\mathcal F\to\mathcal B$，其中 $\mathcal F$ 为 $G$ 的\Term{formal-subgroup}{形式子群}集，$\mathcal B$ 为 $A$ 的同时为\Term{left-coideal}{左余理想}的闭\Term{subalgebra}{子代数}集。于是产生确定其像的问题；下面的命题 5.1 表明，当 $G$ \Term{infinitesimal}{无穷小}时，$Q$ 是\Term{bijection}{双射}。事实上，这对每个 $k$-\Term{formal-group}{形式群} $G$ 都成立。

回顾 2.2.1：\Term{functor}{函子} $G\mapsto\mathrm H(G)$ 是 $k$-\Term{formal-group}{形式群}\Term{category}{范畴}与\Term{cocommutative}{余交换} $k$-\Term{hopf-algebra}{Hopf 代数}\Term{category}{范畴}之间的\Term{equivalence}{等价}。若 $F$ 是由 $A$ 的\Term{closed-ideal}{闭理想} $J$ 定义的 $G$ 的\Term{formal-subgroup}{形式子群}，则 $\mathrm H=\mathrm H(G)$ 的 \Term{Hopf-subalgebra}{Hopf 子代数} $\mathrm H(F)$ 是关于 $A=\mathrm H^*$ 与 $\mathrm H=A^\dagger$ 之间的\Term{dual}{对偶}\Term{pairing}{配对}的 $J$ 的\Term{orthogonality}{正交}\Term{space}{空间}（见 0.2.2）。另一方面，若 $B$ 是同时为\Term{left-coideal}{左余理想}的 $A$ 的闭\Term{subalgebra}{子代数}，则其\Term{orthogonality}{正交}\Term{space}{空间} $I=B^\perp$ 是 $\mathrm H$ 的\Term{coideal}{余理想}（即 $\Delta_{\mathrm H}(I)\subset I\otimes\mathrm H+\mathrm H\otimes I$ 且 $\varepsilon_{\mathrm H}(I)=0$），同时也是\Term{left-ideal}{左理想}。记 $\mathcal H$ 为 $\mathrm H$ 的 \Term{Hopf-subalgebra}{Hopf 子代数}集，$\mathcal I$ 为 $\mathrm H$ 的同时为\Term{coideal}{余理想}的\Term{left-ideal}{左理想}集。对 $I\in\mathcal I$，以 $\pi_I$ 表示\Term{projection}{投影} $\mathrm H\to\mathrm H/I$，$\overline\pi_I$ 表示\Term{projection}{投影} $\mathrm H\to\mathrm H^+\to\mathrm H^+/I$ 的复合，其中 $\mathrm H^+$ 为 $\mathrm H$ 的\Term{augmentation-ideal}{增广理想}。

设 $K$ 为 $\mathrm H$ 的 \Term{Hopf-subalgebra}{Hopf 子代数}，$K^+=K\cap\mathrm H^+$。若 $F$ 是对应于 $K$ 的\Term{formal-subgroup}{形式子群}，则 $J=K^\perp$，而 $A^+/J$ 等同于 $K^+$ 的\Term{dual}{对偶}。又因 $B=Q(F)$ 是映射
\[
A\xrightarrow{\Delta_A}A\widehat\otimes_k A
\xrightarrow{\mathrm{id}\widehat\otimes\overline\pi}A\widehat\otimes_k(A^+/J)
\]
的核，故\Term{dual}{对偶}地，$Q$ 对应于把 $K$ 映为以下复合映射之像的映射 $\Phi$：
\[
\mathrm H\xleftarrow{m_{\mathrm H}}\mathrm H\otimes_k\mathrm H
\xleftarrow{\mathrm{id}\otimes\mathrm{can}.}\mathrm H\otimes_k K^+.
\]
该像即\Term{left-ideal}{左理想} $\mathrm HK^+$，它也是\Term{coideal}{余理想}。同样可见，把 $B\in\mathcal B$ 映为 $\overline{AB^+}$ 的映射\Term{dual}{对偶}地对应于映射 $\Psi$：它把 $I\in\mathcal I$ 映为 $(\mathrm{id}\otimes_k\overline\pi_I)\circ\Delta_{\mathrm H}$ 的核，
\SourcePage{588}
即
\[
\Psi(I)=\{x\in\mathrm H\mid(\mathrm{id}\otimes_k\pi_I)\Delta_{\mathrm H}(x)=x\otimes\pi_I(1)\}.
\]
有以下\Term{theorem}{定理}：
\begin{SGAresult}{\Term{theorem}{定理}}{5.0.1}
设 $k$ 为域，$\mathrm H$ 为\Term{cocommutative}{余交换} \Term{hopf-algebra}{Hopf 代数}。上述映射 $\Phi$ 和 $\Psi$ 是 $\mathrm H$ 的 \Term{Hopf-subalgebra}{Hopf 子代数}集与同时为\Term{coideal}{余理想}的\Term{left-ideal}{左理想}集之间互逆的\Term{bijection}{双射}。
\end{SGAresult}

这一\Term{theorem}{定理}最初由 K.~Newman 证明，见 [Ne75]，\Term{theorem}{定理} 4.1（该处漏写了“\Term{cocommutative}{余交换}”一词）。他的证明利用“Cartier--Gabriel--Kostant \Term{theorem}{定理}”（见 2.9）归结到 $\mathrm H$ \Term{connected}{连通}的情形，再利用这一情形下“Sweedler 基”的存在性（见 [Sw67]，\Term{theorem}{定理} 3）；后者是下文 Dieudonn\'e--Cartier \Term{theorem}{定理} 5.2.2 的\Term{dual}{对偶}结果。H.~J.~Schneider 后来给出一个更短的证明，见 [Sch90]，\Term{theorem}{定理} 4.15。随后 A.~Masuoka 仅假定 $\mathrm H$ 的\Term{coradical}{余根基} $\mathrm H_0$（即所有单\Term{subcoalgebra}{子余代数}之和）可交换，便推广了该结果，见 [Ma91]，\Term{theorem}{定理} 1.3 (3)。

最后指出，对应于\Term{affine}{仿射} $k$-\Term{group-scheme}{群概形} $G$ 的交换 $k$-\Term{hopf-algebra}{Hopf 代数}，若不增加假设，则不能期待 5.0.1 的类似结论，因为对 $G$ 的一个 $k$-\Term{subgroup}{子群} $F$，商 $G/F$ 未必\Term{affine}{仿射}。但 M.~Takeuchi 在 [Tak72]，\Term{theorem}{定理} 4.3（相应地，[Tak79]，\Term{theorem}{定理} 3）中建立了类似的\Term{bijection}{双射}：一端是 $G$ 的\Term{normal}{正规} $k$-\Term{subgroup}{子群} $F$ 的\Term{set}{集合}（相应地，使 $G/F$ \Term{affine}{仿射}的\Term{subgroup}{子群}集），另一端是 $\mathcal O(G)$ 中满足 $\Delta(B)\subset A\otimes B$ 且在\Term{antipode}{对极映射}下\Term{stable}{稳定}的\Term{subalgebra}{子代数} $B$ 的\Term{set}{集合}（相应地，满足 $B\to A$ \Term{faithfully-flat}{忠实平坦}的\Term{subalgebra}{子代数}集）。

\textbf{5.1。} 设 $k$ 为\Term{pseudocompact-ring}{伪紧环}。\footnote{编者注（143）：原文在 5.1 假定 $k$ 为域。事实上可用\Term{flatness}{平坦性}假设替换此假设，因此这里相应修改了 5.1 至 5.1.5。}令 $G$ 为\Term{topologically-flat}{拓扑平坦}的\Term{infinitesimal}{无穷小} $k$-\Term{formal-group}{形式群}，$A$ 为其\Term{affine-algebra}{仿射代数}，$B$ 为 $A$ 的闭\Term{subalgebra}{子代数}，$X=\operatorname{Spf}(B)$，$r:G\to X$ 为包含映射 $B\hookrightarrow A$ 诱导的\Term{epimorphism}{满态射}。我们要研究：在什么条件下，$r$ 使 $X$ 成为 $G$ 模掉一个\Term{subgroup}{子群} $H$ 的右商（见 2.4）。\footnote{编者注（144）：这里把“左”改成了“右”，并修改了命题 5.1 的表述，以更清楚地显示\Term{equivalence}{等价}条件 (i)、(ii) 及结论 $\operatorname{Spf}(B)\simeq G/H$。}

\begin{SGAresult}{命题}{5.1}
设 $G$ 为\Term{topologically-flat}{拓扑平坦}的\Term{infinitesimal}{无穷小} $k$-\Term{formal-group}{形式群}，$A$ 为其\Term{affine-algebra}{仿射代数}，$I_A$ 为 $A$ 的\Term{augmentation-ideal}{增广理想}，$B$ 为 $A$ 的闭\Term{subalgebra}{子代数}，且 $J_B=\overline{AI_B}$，其中 $I_B=B\cap I_A$。假定 $A$ 及对所有 $n\geq0$ 的 $\overline{J_B^n}/\overline{J_B^{n+1}}$ 都在 $k$ 上\Term{topologically-flat}{拓扑平坦}。则以下两个断言\Term{equivalence}{等价}：

(i) 对所有 $x\in B$，$\Delta_A(x)-x\widehat\otimes1$ 属于 $A\widehat\otimes_k J_B$。
\SourcePage{589}
(ii) 下列序列\Term{exact}{正合}，其中 $\tau_1(a)=a\widehat\otimes1$，$\pi$ 为\Term{projection}{投影} $A\to A/J_B$：
\[
\begin{tikzcd}
B\arrow[r]&A\arrow[r,shift left=.7ex,"\tau_1"]\arrow[r,shift right=.7ex,"(\pi\widehat\otimes\mathrm{id}_A)\Delta_A"']&A\widehat\otimes_k(A/J_B).
\end{tikzcd}\tag{*}
\]
即 $B$ 是所有满足 $\Delta_A(x)-x\widehat\otimes1\in A\widehat\otimes_kJ_B$ 的 $x\in A$ 所成的\Term{set}{集合}。

此时，$H=\operatorname{Spf}(A/J_B)$ 是 $G$ 的\Term{formal-subgroup}{形式子群}，且下列序列\Term{exact}{正合}，其中 $\lambda$ 为 $G$ 的\Term{multiplication}{乘法}在 $G\times H$ 上的\Term{restriction}{限制}：
\[
\begin{tikzcd}
G\times H\arrow[r,shift left=.7ex,"\operatorname{pr}_1"]\arrow[r,shift right=.7ex,"\lambda"']&G\arrow[r]&\operatorname{Spf}(B).
\end{tikzcd}\tag{**}
\]
即 $\operatorname{Spf}(B)$ \Term{isomorphism}{同构}于 $G/H$。
\end{SGAresult}
\begin{proof}
置 $\overline A=A/J_B$、$H=\operatorname{Spf}(\overline A)$，则 $H$ 是 $G$ 的\Term{formal-subvariety}{形式子簇}。因 $J_B\subset I_A$，\Term{augmentation}{增广} $\varepsilon_A$ 诱导一个\Term{continuous}{连续} $k$-\Term{algebra}{代数}\Term{morphism}{态射} $\overline\varepsilon:\overline A\to k$。

若 (i) 成立，则 $\Delta_A(I_B)\subset I_B\widehat\otimes1+A\widehat\otimes J_B$，故 $\Delta_A$ 通过取商诱导一个\Term{diagonal-morphism}{对角态射} $\overline\Delta$。于是 $\overline\Delta$ 和 $\overline\varepsilon$ 赋予 $H$ 一个 $G$ 的\Term{formal-submonoid}{形式子幺半群}结构。因 $G$ \Term{infinitesimal}{无穷小}，$H$ 也\Term{infinitesimal}{无穷小}；由命题 2.7，$H$ 是 $G$ 的\Term{formal-subgroup}{形式子群}。由 $G/H$ 的定义（见 2.4），(ii) \Term{pushout}{推出}命题的最后一个断言。

另一方面，显然 (ii) \Term{pushout}{推出} (i)。逆命题的证明占据 5.1.1 至 5.1.5。
\end{proof}

\textbf{5.1.1。} 先考虑如下\Term{category}{范畴} $\mathcal C$：其\Term{object}{对象}为偶对 $(A,J)$，其中 $A$ 是一个\Term{profinite-algebra}{预有限} $k$-\Term{algebra}{代数}，$J$ 是 $A$ 的\Term{closed-ideal}{闭理想}；$\mathcal C$ 的\Term{morphism}{态射} $\psi:(A,J)\to(A',J')$ 是把 $J$ 映入 $J'$ 的\Term{continuous}{连续} $k$-\Term{algebra}{代数}\Term{homomorphism}{同态} $A\to A'$。把 $(A,J)$ 映为 $(\operatorname{Spf}(A/J),\operatorname{Spf}(A))$，显然得到 $\mathcal C$ 到如下偶对 $(Z,Y)$ 的\Term{category}{范畴}的\Term{anti-equivalence}{反等价}：$Y$ 为 $k$-\Term{formal-variety}{形式簇}，$Z$ 为其\Term{formal-subvariety}{形式子簇}；\Term{morphism}{态射} $\phi:(Z,Y)\to(Z',Y')$ 是把 $Z$ 映入 $Z'$ 的 $k$-\Term{formal-variety}{形式簇}\Term{morphism}{态射} $Y\to Y'$。

$\mathcal C$ 的一个\Term{object}{对象} $(A,J)$ 上的\Term{cogroup-object}{余群}结构，是 $\operatorname{Spf}(A)$ 上满足以下条件的一个\Term{formal-group}{形式群}结构（沿用 2.1 的记号）：
\[
\begin{aligned}
{}&(1)\quad\Delta_A(J)\subset J\widehat\otimes_k A+A\widehat\otimes_kJ;\\
&(2)\quad\varepsilon_A(J)=0;\\
&(3)\quad c_A(J)\subset J.
\end{aligned}
\]
这些条件也就是说，$H=\operatorname{Spf}(A/J)$ 是 $G=\operatorname{Spf}(A)$ 的\Term{formal-subgroup}{形式子群}。\footnote{编者注（145）：若 $(A',J')$ 是 $\mathcal C$ 的第二个\Term{cogroup-object}{余群}，对应于 $k$-\Term{formal-group}{形式群}偶对 $H'\subset G'$，则给定\Term{cogroup-object}{余群}\Term{morphism}{态射} $(A',J')\to(A,J)$ \Term{equivalence}{等价}于给定把 $H$ 映入 $H'$ 的 $k$-\Term{formal-group}{形式群}\Term{morphism}{态射} $G\to G'$。}
\SourcePage{590}
再设 $A$ 为\Term{local-ring}{局部环}，即 $\operatorname{Spf}(A)$ 为\Term{infinitesimal}{无穷小}\Term{formal-group}{形式群}。则当 $J\ne A$ 时，条件 (2)、(3) 是 (1) 的推论。事实上，若 $J$ 是异于 $A$ 的\Term{closed-ideal}{闭理想}，则 $J\subset I_A$，所以 (2) 成立，而 $M=\operatorname{Spf}(A/J)$ 是 $G$ 的\Term{formal-submonoid}{形式子幺半群}。因 $G$ \Term{infinitesimal}{无穷小}，由 2.7，$M$ 是 $G$ 的\Term{formal-subgroup}{形式子群}，即条件 (3) 成立。

\textbf{5.1.2。} 以 $\mathbf{Alpg}/k$ 表示“\Term{graded}{分次}\Term{profinite-algebra}{预有限}”$k$-\Term{algebra}{代数}\Term{category}{范畴}：其\Term{object}{对象}为一列\Term{pseudocompact}{伪紧} $k$-模 $A_0,A_1,\ldots,A_n,\ldots$ 及乘积 $\prod_{n\geq0}A_n$ 上的一个\Term{profinite-algebra}{预有限}\Term{algebraic-structure}{代数结构}，满足 $A_nA_m\subset A_{m+n}$（借\Term{canonical}{典范}\Term{injectivity}{单射}将 $A_n$ 等同于 $\prod_{i\geq0}A_i$ 的\Term{subset}{子集}）；\Term{morphism}{态射} $\psi:(A_n)\to(B_n)$ 为一列\Term{continuous}{连续}\Term{linear}{线性}映射 $\psi_n:A_n\to B_n$，满足：若 $a\in A_m$、$a'\in A_n$，则 $\psi_{m+n}(aa')=\psi_m(a)\psi_n(a')$。

\textbf{定义。} 显然，两个\Term{graded}{分次}\Term{profinite-algebra}{预有限} $k$-\Term{algebra}{代数} $(A_n)$ 和 $(B_n)$ 在 $\mathbf{Alpg}/k$ 中有\Term{coproduct}{余积}，\footnote{编者注（146）：这里把“\Term{direct-sum}{直和}”改成了“\Term{coproduct}{余积}”。}其第 $n$ 个\Term{set-component}{分量}为\Term{pseudocompact}{伪紧} $k$-模 $A_i\widehat\otimes_k B_{n-i}$ 的乘积 $\prod_{i=0}^n A_i\widehat\otimes_k B_{n-i}$。把此\Term{coproduct}{余积}记为 $(A_n)\widehat\otimes_k(B_n)$。

于是，$\mathbf{Alpg}/k$ 的\Term{object}{对象} $(A_n)$ 上的\Term{cogroup-object}{余群}结构，由\Term{continuous}{连续} $k$-\Term{linear}{线性}映射
\[
\Delta_n:A_n\longrightarrow\prod_{i=0}^n A_i\widehat\otimes_k A_{n-i},\qquad
\varepsilon:A_0\longrightarrow k
\]
给出，它们在 $\prod_{n\geq0}A_n$ 上诱导 $\mathbf{Alp}/k$ 中的\Term{cogroup-object}{余群}结构（对 $i\geq1$ 置 $\varepsilon(A_i)=0$）。

最后，对 $\mathcal C$ 的每个\Term{object}{对象} $(A,J)$，记 $\operatorname{Gr}_J(A)$ 为 $A$ 关于 $J$ 的幂的\Term{closure}{闭包} $\overline{J^n}$ 所定义\Term{filtration}{滤过}的\Term{associated-graded}{关联分次}\Term{profinite-algebra}{预有限}\Term{algebra}{代数}。因此 $\operatorname{Gr}_J(A)_n=\overline{J^n}/\overline{J^{n+1}}$，其\Term{multiplication}{乘法}由 $A$ 的\Term{multiplication}{乘法}诱导。

\begin{SGAresult}{\Term{lemma}{引理}}{5.1.2}
设 $U,V$ 为两个\Term{pseudocompact}{伪紧} $k$-模，$U$ \Term{topologically-flat}{拓扑平坦}；设 $U=U_0\supset U_1\supset\cdots$ 及 $V=V_0\supset V_1\supset\cdots$ 为两个由闭 $k$-\Term{submodule}{子模}组成的递降列。\footnote{编者注（147）：原文仅在 $k$ 为域时陈述此\Term{lemma}{引理}，并把该情形的证明留给读者。}用闭\Term{submodule}{子模}
\[
W_n=U_n\widehat\otimes_kV_0+U_{n-1}\widehat\otimes_kV_1+\cdots+U_0\widehat\otimes_kV_n
\]
\Term{filtration}{滤过}\Term{completed-tensor-product}{完备张量积} $W=U\widehat\otimes_kV$。假定每个 $U_i/U_{i+1}$ 在 $k$ 上\Term{topologically-flat}{拓扑平坦}（从而对每个 $n$，$U/U_n$ 及 $U_n$ 也\Term{topologically-flat}{拓扑平坦}）。则对每个 $n$，有\Term{isomorphism}{同构}
\[W_n/W_{n+1}\simeq\bigoplus_{i+j=n}(U_i/U_{i+1})\widehat\otimes_k(V_j/V_{j+1}).\]
\end{SGAresult}
\begin{proof}
对 $i,j\geq0$，置 $W_{i,j}=U_i\widehat\otimes V_j$、$\overline W_{i,j}=(U_i/U_{i+1})\widehat\otimes_k(V_j/V_{j+1})$。对 $n$ 归纳，证明自然映射
\[\pi_n:W_n\longrightarrow\bigoplus_{i+j=n}\overline W_{i,j}\]
\SourcePage{591}
为\Term{surjectivity}{满射}，且包含关系 $W_{n+1}\subset\operatorname{Ker}(\pi_n)$ 为等号。当 $n=0$ 时，\Term{projection}{投影}
\[\pi_0:U_0\widehat\otimes_kV_0\longrightarrow(U_0/U_1)\widehat\otimes_k(V_0/V_1)\]
为\Term{surjectivity}{满射}；因 $U_0$、$U_0/U_1$ 及 $U_1$ 在 $k$ 上\Term{topologically-flat}{拓扑平坦}，可见
\[
\operatorname{Ker}(\pi_0)=U_0\widehat\otimes_kV_1+U_1\widehat\otimes_kV_0,\qquad
(U_0\widehat\otimes_kV_1)\cap(U_1\widehat\otimes_kV_0)=U_1\widehat\otimes_kV_1.
\]
现设 $n>0$，且结果对 $n-1$ 成立。置 $M_0=U_0\widehat\otimes_kV_n$、$S_0=\sum_{i=1}^nU_i\widehat\otimes_kV_{n-i}$。有 $S_0\subset U_1\widehat\otimes_kV_0$，故将上面的论证用于 $V_0\supset V_n$ 而非 $V_0\supset V_1$，得到
\[M_0\cap S_0\subset(U_0\widehat\otimes_kV_n)\cap(U_1\widehat\otimes_kV)=U_1\widehat\otimes_kV_n.\]
由此得到 $M_0\cap S_0=U_1\widehat\otimes_kV_n$。因 $W_n=M_0+S_0$，得到下列行\Term{exact}{正合}的\Term{commutative-diagram}{交换图}，其中对 $i=0,\ldots,n-1$，置 $U'_i=U_{i+1}$、$W'_{i,n-1-i}=W_{i+1,n-1-i}$：
\[
\begin{tikzcd}[column sep=small]
0\arrow[r]&S_0\arrow[r]\arrow[d,"\pi'_{n-1}"']&W_n\arrow[r]\arrow[d,"\pi_n"]&(U_0/U_1)\widehat\otimes_kV_n\arrow[r]\arrow[d,"p"]&0\\
0\arrow[r]&\bigoplus_{i=0}^{n-1}\overline W'_{i,n-1-i}\arrow[r]&\bigoplus_{i=0}^n\overline W_{i,n-1}\arrow[r]&\overline W_{0,n}\arrow[r]&0.
\end{tikzcd}
\]
则 $p$ 为\Term{surjectivity}{满射}，核为 $(U_0/U_1)\widehat\otimes_kV_{n+1}$。又将归纳假设用于列 $(U'_i)$，可知 $\pi'_{n-1}$ 为\Term{surjectivity}{满射}，其核为 $W'_n=\sum_{i=1}^nW_{i,n+1-i}$。于是 $\pi_n$ 为\Term{surjectivity}{满射}，且 $W_{n+1}\subset\operatorname{Ker}(\pi_n)$ 为等号。\Term{lemma}{引理}得证。
\end{proof}

回到 $\mathcal C$ 的\Term{object}{对象} $(A,J)$。注意由 0.2.G，每个 $\overline{J^n}/\overline{J^{n+1}}$ 在 $k$ 上\Term{topologically-flat}{拓扑平坦}，\Term{equivalence}{等价}于 $\operatorname{Gr}_J(A)$ 在 $k$ 上\Term{topologically-flat}{拓扑平坦}。
\begin{SGAresult}{推论}{5.1.2}
设 $\mathcal P$ 为 $\mathcal C$ 的如下\Term{full-subcategory}{满子范畴}：其\Term{object}{对象} $(A,J)$ 满足 $A$ 和 $\operatorname{Gr}_J(A)$ 均在 $k$ 上\Term{topologically-flat}{拓扑平坦}。则\Term{functor}{函子} $\mathcal P\to\mathbf{Alpg}/k$，$(A,J)\mapsto\operatorname{Gr}_J(A)$，与\Term{finite}{有限}\Term{coproduct}{余积}可交换，因而把 $\mathcal P$ 的\Term{cogroup-object}{余群}变为 $\mathbf{Alpg}/k$ 的\Term{cogroup-object}{余群}。

特别地，若 $k$ 为域，则对 $\mathcal C$ 的每个\Term{cogroup-object}{余群} $(A,J)$，$\operatorname{Gr}_J(A)$ 是 $\mathbf{Alpg}/k$ 的\Term{cogroup-object}{余群}。\footnote{编者注（148）：因此，对 $k$-\Term{formal-group}{形式群}偶对 $H\subset G$，得到“$G$ 沿 $H$ 的\Term{formal-completion}{形式完备化} $\widehat G_H$”，它是一个 $k$-\Term{formal-group}{形式群}；5.1.3--5.1.4 还将表明，包含映射 $\sigma:H\hookrightarrow\widehat G_H$ 有一个\Term{retraction}{收缩} $\pi:\widehat G_H\to H$，而 $k$-\Term{formal-group}{形式群} $N=\operatorname{Ker}(\pi)$ 作为\Term{formal-variety}{形式簇}等同于\Term{homogeneous-space}{齐性空间} $G/H$ 沿\Term{unit-section}{单位截面}的\Term{completion}{完备化}。这将用于 5.2.2。}
\end{SGAresult}

\textbf{5.1.3。} 将每个\Term{profinite-algebra}{预有限} $k$-\Term{algebra}{代数} $\Gamma$ 等同于满足 $\Gamma_0=\Gamma$、$n>0$ 时 $\Gamma_n=0$ 的\Term{graded}{分次}\Term{profinite-algebra}{预有限} $k$-\Term{algebra}{代数} $(\Gamma_n)_{n\geq0}$。特别地，若 $(A_n)_{n\geq0}$ 为\Term{graded}{分次}\Term{profinite-algebra}{预有限} $k$-\Term{algebra}{代数}，可同时把 $A_0$ 看作\Term{profinite-algebra}{预有限} $k$-\Term{algebra}{代数}和\Term{graded}{分次}\Term{profinite-algebra}{预有限} $k$-\Term{algebra}{代数}。以 $\rho:(A_n)\to A_0$ 表示 $\mathbf{Alpg}/k$ 中满足 $\rho_0=\mathrm{id}_{A_0}$、$n>0$ 时 $\rho_n=0$ 的\Term{morphism}{态射}。
\SourcePage{592}
同样，以 $\tau:A_0\to(A_n)$ 表示满足 $\tau_0=\mathrm{id}_{A_0}$、$n>0$ 时 $\tau_n=0$ 的 $\rho$ 的\Term{section}{截面}。

$(A_n)_{n\in\mathbf N}$ 上的每个\Term{cogroup-object}{余群}结构在 $A_0$ 上诱导一个\Term{cogroup-object}{余群}结构，使 $\rho$、$\tau$ 为\Term{cogroup-object}{余群}\Term{homomorphism}{同态}。此时，记 $I_0$ 为 $A_0$ 的\Term{augmentation-ideal}{增广理想}，对 $n\geq0$ 置 $A'_n=A_n/\overline{I_0A_n}$（从而 $A'_0=k$）。

于是 $(A'_n)_{n\in\mathbf N}$ 为 $\mathbf{Alpg}/k$ 的\Term{cogroup-object}{余群}（注意因 $A_0=I_0\oplus k\cdot1$，$I_0\widehat\otimes_kA_n\simeq\overline{I_0A_n}$ 是 $A_n$ 的\Term{direct-summand}{直和因子}）。因 $\tau$ 为 $\rho$ 的\Term{section}{截面}，由 2.4.A，\Term{cogroup-object}{余群} $(A_n)_{n\in\mathbf N}$ 是 $A_0$ 与\Term{cogroup-object}{余群} $(A'_n)_{n\in\mathbf N}$ 的“\Term{semidirect-coproduct}{半直余积}”。准确地说，$(A'_n)$ 作为 $\mathbf{Alpg}/k$ 的\Term{object}{对象}，\Term{isomorphism}{同构}于以下映射对的核：
\[
\begin{tikzcd}(A_n)\arrow[r,shift left=.7ex,"\tau_1"]\arrow[r,shift right=.7ex,"(\mathrm{id}\widehat\otimes\rho)\Delta"']&(A_n)\widehat\otimes_kA_0.\end{tikzcd}
\]
其中 $\Delta:(A_n)\to(A_n)\widehat\otimes(A_n)$ 为\Term{comultiplication}{余乘法}，$\tau_1(x)=x\widehat\otimes1$。将 $A'_n$ 等同于它在 $A_n$ 中的像，则映射
\[(A'_n\widehat\otimes_kA_0)\longrightarrow(A_n),\qquad a'_n\widehat\otimes a_0\longmapsto a'_na_0\]
是 $\mathbf{Alpg}/k$ 的\Term{isomorphism}{同构}。（注意：它并非\Term{cogroup-object}{余群}\Term{isomorphism}{同构}，但 $\Delta(A')\subset A\widehat\otimes A'$，且 $(\gamma\rho'\widehat\otimes\mathrm{id})\circ\Delta|_{A'}=\Delta'$，其中 $\Delta'$ 为 $A'$ 的\Term{comultiplication}{余乘法}，$\gamma\rho'$ 为\Term{projection}{投影} $A\to A'$，见 2.4.A。）

\textbf{5.1.4。} 设 $(A,J)$ 为 $\mathcal C$ 的\Term{object}{对象}，$(A_n)=\operatorname{Gr}_J(A)$ 为对应的 $\mathbf{Alpg}/k$ 的\Term{object}{对象}，即对 $n\geq0$，$A_n=\overline{J^n}/\overline{J^{n+1}}$。显然\Term{algebra}{代数} $\mathscr A=\prod_{n\geq0}A_n$ 由 $A_0$ 和 $A_1$ 生成，即当 $n\geq1$ 时，由\Term{multiplication}{乘法}定义的映射
\[
A_1\widehat\otimes_{A_0}A_1\widehat\otimes_{A_0}\cdots\widehat\otimes_{A_0}A_1\longrightarrow A_n\tag{1}
\]
是\Term{surjectivity}{满射}。

再设 $(A,J)$ 为 $\mathcal C$ 的\Term{cogroup-object}{余群}，且 $A$ 及商 $\overline{J^n}/\overline{J^{n+1}}$ 在 $k$ 上\Term{flat}{平坦}。由推论 5.1.2，$\operatorname{Gr}_J(A)$ 为 $\mathbf{Alpg}/k$ 的\Term{cogroup-object}{余群}。所以由 5.1.3，若置
\[
A'_n=\left\{x\in\overline{J^n}/\overline{J^{n+1}}\;\middle|\;
\Delta(x)-x\widehat\otimes1\in\bigoplus_{i=1}^n
(\overline{J^{n-i}}/\overline{J^{n-i+1}})\widehat\otimes
(\overline{J^i}/\overline{J^{i+1}})\right\},\tag{2}
\]
则映射 $(A'_n\widehat\otimes A_0)\to(A_n)$，$a'_n\widehat\otimes a_0\mapsto a'_na_0$，为 $\mathbf{Alpg}/k$ 的\Term{isomorphism}{同构}。\footnote{编者注（149）：设 $G$、$H$、$\widehat G_H$ 分别是对应于 $A$、$A_0=A/J$、$\operatorname{Gr}_J(A)$ 的 $k$-\Term{formal-group}{形式群}。则 $\tau:A_0\hookrightarrow\operatorname{Gr}_J(A)$ 对应于包含映射 $H\hookrightarrow\widehat G_H$ 的\Term{retraction}{收缩} $\pi:\widehat G_H\to H$；以上结论意味着 $\widehat G_H$ 为 $N=\operatorname{Ker}(\pi)$ 与 $H$ 的\Term{semidirect-product}{半直积}。}
\SourcePage{593}
记 $I_0$ 为 $A_0$ 的\Term{augmentation-ideal}{增广理想}。由 (1) 及以下\Term{commutative-diagram}{交换图}（其中 $A_1'{}^{\widehat\otimes n}=A'_1\widehat\otimes_k\cdots\widehat\otimes_kA'_1$，共 $n$ 个因子），
\[
\begin{tikzcd}[column sep=small]
A_1\widehat\otimes_{A_0}\cdots\widehat\otimes_{A_0}A_1\arrow[d,"m"']&A_1'{}^{\widehat\otimes n}\widehat\otimes_k A_0\arrow[l,"\sim"']\arrow[d,"m'\widehat\otimes\mathrm{id}"]&(A_1'{}^{\widehat\otimes n}\widehat\otimes_k I_0)\oplus A_1'{}^{\widehat\otimes n}\arrow[l,"\sim"']\arrow[d,"{m'\widehat\otimes\mathrm{id}\oplus m'}"]\\
A_n&A'_n\widehat\otimes_kA_0\arrow[l,"\sim"']&(A'_n\widehat\otimes_kI_0)\oplus A'_n\arrow[l,"\sim"']
\end{tikzcd}
\]
可知由\Term{multiplication}{乘法}诱导的映射
\[m':A'_1\widehat\otimes_k\cdots\widehat\otimes_kA'_1\longrightarrow A'_n\tag{3}\]
为\Term{surjectivity}{满射}；换言之，\Term{profinite-algebra}{预有限} $k$-\Term{algebra}{代数} $\mathscr A'=\prod_{n\geq0}A'_n$ 由其 $1$ 次项生成。

现回到命题 5.1 的假设 (i)：\footnote{编者注（150）：以下详述了原文，并把“补充”$\overline{I_B^n}=B\cap\overline{J_B^n}$ 移到最后（命题 5.1 的证明不需要此补充）。}设 $G$ 为\Term{topologically-flat}{拓扑平坦}的\Term{infinitesimal}{无穷小} $k$-\Term{formal-group}{形式群}，$A$ 为其\Term{affine-algebra}{仿射代数}，$I_A$ 为\Term{augmentation-ideal}{增广理想}，$B$ 为 $A$ 的闭\Term{subalgebra}{子代数}，$J=\overline{AI_B}$，其中 $I_B=B\cap I_A$。记 $H$ 为\Term{formal-subgroup}{形式子群} $\operatorname{Spf}(A/J)$，$\pi$ 为\Term{projection}{投影} $A\to A/J$。设对每个 $n\geq1$，$A/\overline{J^n}$ 在 $k$ 上\Term{topologically-flat}{拓扑平坦}，且 $B$ 包含于以下映射对的核 $\widetilde B$：
\[
\begin{tikzcd}A\arrow[r,shift left=.7ex,"\tau_1"]\arrow[r,shift right=.7ex,"(\mathrm{id}_A\widehat\otimes_k\pi)\Delta_A"']&A\widehat\otimes_k(A/J).\end{tikzcd}
\]
设 $(A_n)=\operatorname{Gr}_J(A)$，$I_0=I_A/J$ 为 $A_0=A/J$ 的\Term{augmentation-ideal}{增广理想}，并如上方 (2) 定义 $(A'_n)_{n\in\mathbf N}$。记 $(B_n)_{n\in\mathbf N}$、$(\widetilde B_n)_{n\in\mathbf N}$ 为 $B$、$\widetilde B$ 上由 $A$ 的\Term{filtration}{滤过}诱导的\Term{filtration}{滤过}所对应的 $\mathbf{Alpg}/k$ 的\Term{object}{对象}；这些\Term{filtration}{滤过}分别由\Term{ideal}{理想} $B\cap\overline{J^n}$、$\widetilde B\cap\overline{J^n}$ 定义。显然，对每个 $n$，$B_n\subset\widetilde B_n\subset A'_n$，而
\[\mathscr B=\prod_{n\geq0}B_n\subset\widetilde{\mathscr B}=\prod_{n\geq0}\widetilde B_n\]
是 $\mathscr A'=\prod_{n\geq0}A'_n$ 的\Term{subalgebra}{子代数}。

另一方面，$J$（相应地，$\overline{J^2}$）为 $I_B\widehat\otimes_kA$（相应地，$\overline{I_B^2}\widehat\otimes_k A$）在 $A$ 中的像。因此，映射
\[(I_B/\overline{I_B^2})\widehat\otimes_k A\longrightarrow J/\overline{J^2}=A_1\simeq A'_1\widehat\otimes_kA_0\]
\SourcePage{594}
为\Term{surjectivity}{满射}，且经由 $(I_B/\overline{I_B^2})\widehat\otimes_kA_0$ \Term{resolution}{分解}。因 $A_0=k\cdot1\oplus I_0$，由下列\Term{exact}{正合}\Term{commutative-diagram}{交换图}
\[
\begin{tikzcd}
(I_B/\overline{I_B^2})\widehat\otimes_kA_0\arrow[d]&(I_B/\overline{I_B^2})\widehat\otimes_kI_0\oplus(I_B/\overline{I_B^2})\arrow[l,"\sim"']\arrow[d]\\
J/\overline{J^2}\arrow[d]&A'_1\widehat\otimes_kI_0\oplus A'_1\arrow[l,"\sim"']\arrow[d]\\
0&0
\end{tikzcd}
\]
可知 $A'_1$ 为 $I_B/\overline{I_B^2}$ 的像，故 $B_1=A'_1$。因 $\mathscr A'$ 由 $A'_1$ 生成，对所有 $n\geq1$ 有
\[A'_n\subset B_n\subset\widetilde B_n\subset A'_n,\]
所以 $B_n=\widetilde B_n=A'_n$。\footnote{编者注（151）：沿用编者注（149）的记号，这说明 $G/H$ 沿\Term{unit-section}{单位截面}的\Term{formal-completion}{形式完备化}（其\Term{affine-algebra}{仿射代数}为 $\prod_nB_n$）作为\Term{formal-variety}{形式簇}\Term{isomorphism}{同构}于 $k$-\Term{formal-group}{形式群} $N$。}

最后，因\Term{formal-group}{形式群} $G$ \Term{infinitesimal}{无穷小}，有 $\mathfrak r(A)=I_A$，故\Term{ideal}{理想} $\overline{I_A^n}$ 趋于 $0$（见 0.1.2）；于是 $\overline{J^n}$ 也趋于 $0$，从而 $\widetilde B$ 和 $B$ 上的诱导\Term{filtration}{滤过}均\Term{hausdorff}{分离}。又因 $B$ 为 $A$ 的闭\Term{subalgebra}{子代数}，它关于\Term{ideal}{理想} $B\cap\overline{J^n}$ 所定义的\Term{topology}{拓扑}\Term{complete}{完备}。由 [CA]，第 V 节 7，\Term{lemma}{引理} 1（另见下面的 5.1.5），$B=\widetilde B$。命题 5.1 的证明至此完成。

还有如下补充。对每个 $n$，$\overline{J^n}=\overline{AI_B^n}$ 是 $A\widehat\otimes_B\overline{I_B^n}$ 在 $A$ 中的像，也是 $A\widehat\otimes_B(B\cap\overline{J^n})$ 的像。由假设，\Term{formal-subgroup}{形式子群} $H$ 的\Term{affine-algebra}{仿射代数} $A/J$ 在 $k$ 上\Term{topologically-flat}{拓扑平坦}。因此，由\Term{theorem}{定理} 2.4，\Term{morphism}{态射} $G=\operatorname{Spf}(A)\to G/H=\operatorname{Spf}(B)$ 是满的且\Term{topologically-flat}{拓扑平坦}。于是
\[A\widehat\otimes_B\overline{I_B^n}=\overline{J^n}=A\widehat\otimes_B(B\cap\overline{J^n}),\]
这\Term{pushout}{推出}对每个 $n$，$\overline{I_B^n}=B\cap\overline{J^n}$。也可由映射
\[\overline{I_B^n}/\overline{I_B^{n+1}}\longrightarrow A'_i=(B\cap\overline{J^n})/(B\cap\overline{J^{n+1}})\]
为\Term{surjectivity}{满射}，及下面的 5.1.5 (ii)（取 $B'_n=\overline{I_B^n}$、$B_n=B\cap\overline{J^n}$）得到。

\begin{SGAresult}{\Term{lemma}{引理}}{5.1.5}
(i) 设 $M,N$ 为两个\Term{abelian-group}{阿贝尔群}，分别由递降\Term{subgroup}{子群}列 $(M_n)_{n\in\mathbf Z}$、$(N_n)_{n\in\mathbf Z}$ \Term{filtration}{滤过}。假定 $\bigcup_nM_n=M$、$\bigcup_nN_n=N$，$\bigcap_nM_n=\bigcap_nN_n=0$，且 $M$ 关于 $M_n$ 定义的\Term{topology}{拓扑}\Term{complete}{完备}。设 $f:M\to N$ 为\Term{filtration}{滤过}群\Term{morphism}{态射}。\footnote{编者注（152）：原文仅陈述 (ii)；为方便读者，这里在 (i) 中陈述了 [CA]，第 V 节 7 的\Term{lemma}{引理} 1。}
\SourcePage{595}

(a) 若 $f$ 在\Term{associated-graded}{关联分次}上诱导\Term{surjectivity}{满射}，则 $f$ 是\Term{surjectivity}{满射}，且 $N$ 关于 $N_n$ 定义的\Term{topology}{拓扑}\Term{complete}{完备}。

(b) 若 $f$ 在\Term{associated-graded}{关联分次}上诱导\Term{injectivity}{单射}，则 $f$ 是\Term{injectivity}{单射}。

(ii) 设 $B$ 为\Term{abelian-group}{阿贝尔群}，$B=B'_0\supset B'_1\supset\cdots$ 和 $B=B_0\supset B_1\supset\cdots$ 为 $B$ 的两个\Term{hausdorff}{分离}\Term{subgroup}{子群}\Term{filtration}{滤过}，且对每个 $n$，$B'_n\subset B_n$。设 $B$ 关于\Term{filtration}{滤过} $(B'_n)$ 定义的\Term{topology}{拓扑}\Term{complete}{完备}。

若对每个 $i$，映射 $B'_i/B'_{i+1}\to B_i/B_{i+1}$ 为\Term{surjectivity}{满射}，则对每个 $n$，$B'_n=B_n$。
\end{SGAresult}
\begin{proof}
(i) 是 [CA]，第 V 节 7 的\Term{lemma}{引理} 1（另见 [BAC]，III，第 2.8 节）；(ii) 由 (i) 得到，只需取 $M=B'_n\supset B'_{n+1}\supset\cdots$ 及 $N=B_n\supset B_{n+1}\supset\cdots$。
\end{proof}

\textbf{5.2。} 在第 5 节以下各段中，$k$ 表示\Term{characteristic}{特征} $p>0$ 的\Term{perfect-field}{完美域}。

置 $\overline{\mathbf N}=\mathbf N\cup\{\infty\}$。若 $B$ 为\Term{profinite-algebra}{预有限} $k$-\Term{algebra}{代数}，$r\in\mathbf N$，则以 $((x^{p^r}))_{x\in\mathfrak r(B)}$ 表示由 $B$ 的\Term{radical}{根基} $\mathfrak r(B)$ 中元素 $x$ 的 $p^r$ 次幂生成的 $B$ 的\Term{closed-ideal}{闭理想}。若 $r=\infty$，仍使用这一记号，约定 $((x^{p^\infty}))_{x\in\mathfrak r(B)}$ 为零\Term{ideal}{理想}。在这两种情形，$B_r$ 均表示商 $B/((x^{p^r}))_{x\in\mathfrak r(B)}$。

若 $((x^{p^r}))_{x\in\mathfrak r(B)}$ 为零\Term{ideal}{理想}，则称 $B$ 的\Term{height}{高度} $\leq r$；若此时 $r$ \Term{finite}{有限}，则称 $B$ 为\Term{finite}{有限}\Term{height}{高度}的。

特别考虑 $B=k[[\omega]]$ 的情形，其中 $\omega$ 为\Term{pseudocompact}{伪紧} $k$-\Term{vector-space}{向量空间}（见 1.2.5）。\footnote{编者注（153）：原文使用“\Term{linearly-compact-vector-space}{线性紧向量空间}”，它\Term{equivalence}{等价}于“\Term{pseudocompact}{伪紧}\Term{vector-space}{向量空间}”（见 [BAC]，III，第 2 节，习题 15 a)、19 a)、20 d)）。这里沿用此前使用的“\Term{pseudocompact}{伪紧}”一词。}此时称 $B$ 为\Term{formal-power-series-algebra}{形式幂级数代数}，$B_r$ 为\Term{truncated-formal-power-series-algebra}{截断形式幂级数代数}（$r\in\overline{\mathbf N}$；因此约定 $B=B_\infty$ 也称为“\Term{truncation}{截断}的”）。若 $B=k[[\omega]]$，也把 $((x^{p^r}))_{x\in\mathfrak r(B)}$ 写作 $((x^{p^r}))_{x\in\omega}$。

\textbf{记号。} 设 $\omega$ 为\Term{pseudocompact}{伪紧} $k$-\Term{vector-space}{向量空间}，由一列递增闭向量子\Term{space}{空间}\Term{filtration}{滤过}：
\[0=\omega_0\subset\omega_1\subset\omega_2\subset\omega_3\subset\cdots.\]
(a) $k[[\omega]]$ 中由所有 $x^{p^r}$（$r\in\mathbf N$，$x\in\omega_r$）生成的\Term{closed-ideal}{闭理想}记为 $((x^{p^r}))_{r\in\omega_r}$。

(b) 另一方面，以 ${}_r\omega$ 表示满足下式的\Term{filtration}{滤过}\Term{pseudocompact}{伪紧}\Term{vector-space}{向量空间}：
\[{}_r\omega_i=\omega_i\quad(i<r),\qquad{}_r\omega_i=\omega\quad(i\geq r).\]

\begin{SGAresult}{\Term{theorem}{定理}（Dieudonn\'e--Cartier）}{5.2}
设 $H\to G$ 为\Term{characteristic}{特征} $p>0$ 的\Term{perfect-field}{完美域} $k$ 上\Term{infinitesimal}{无穷小}\Term{formal-group}{形式群}之间的一个\Term{monomorphism}{单态射}。设 $B$ 为\Term{homogeneous-space}{齐性空间} $G/H$ 的\Term{affine-algebra}{仿射代数}，并假定以下三个条件之一成立：\footnote{编者注（154）：这里将 $H\backslash G$ 改成 $G/H$，证明中亦同；并增加了条件 (iii)。}
\SourcePage{596}

(i) $B$ 为\Term{finite}{有限}\Term{height}{高度}的（特别地，若 $G$ 为\Term{finite}{有限}\Term{height}{高度}的，则此条件成立）。

(ii) $B$ 为\Term{complete-noetherian-local-ring}{完备诺特局部环}。

(iii) $B$ 为\Term{reduced}{约化}环。

则 $B$ \Term{isomorphism}{同构}于\Term{finite}{有限}族\Term{truncated-formal-power-series-algebra}{截断形式幂级数代数}的\Term{completed-tensor-product}{完备张量积}。
\end{SGAresult}
本\Term{theorem}{定理}的证明占据 5.2.1 至 5.2.5。

\textbf{5.2.1。} 设 $A$ 为 $G$ 的\Term{affine-algebra}{仿射代数}，$I_A$ 为其\Term{augmentation-ideal}{增广理想}，$I=B\cap I_A$。由 2.4，$H=\operatorname{Spf}(A/\overline{AI})$，且
\[B=\{x\in A\mid\Delta(x)-1\widehat\otimes x\in A\widehat\otimes\overline{AI}\}.\]
置 $\omega=I/\overline{I^2}$。以 $I_r$ 表示 $I$ 中由满足 $x^{p^r}=0$ 的元素组成的闭子\Term{ideal}{理想}，$\omega_r$ 表示 $I_r$ 在 $\omega$ 中的\Term{canonical}{典范}像。我们要证明：
\begin{SGAresult}{命题}{5.2.1}
若\Term{projection}{投影} $I\to I/\overline{I^2}$ 存在\Term{continuous}{连续}\Term{section}{截面} $\sigma:\omega\to I$，对每个 $r$ 均满足 $\sigma(\omega_r)\subset I_r$，则 $B$ \Term{isomorphism}{同构}于 $k[[\omega]]/((x^{p^r}))_{x\in\omega_r}$。
\end{SGAresult}
这样的\Term{section}{截面}延拓为一个\Term{continuous}{连续}\Term{morphism}{态射} $k[[\omega]]\to B$，且经由 $B'=k[[\omega]]/((x^{p^r}))_{x\in\omega_r}$ \Term{resolution}{分解}。我们将在 5.2.2 至 5.2.5 证明由此得到的\Term{morphism}{态射}
\[\phi:B'=k[[\omega]]/((x^{p^r}))_{x\in\omega_r}\longrightarrow B\]
是\Term{isomorphism}{同构}。

\textbf{5.2.1.A。}\footnote{编者注（155）：5.2.1.A 和 5.2.1.B 为编者所加。}对每个 $r\in\mathbf N$，置 $\mathcal I_r=\overline{I^2}+I_r$，从而 $\mathcal I_r/\overline{I^2}\simeq\omega_r$；有行\Term{exact}{正合}的\Term{commutative-diagram}{交换图}
\[
\begin{tikzcd}
0\arrow[r]&\mathcal I_{r-1}\arrow[r]\arrow[d]&\mathcal I_r\arrow[r]\arrow[d]&\mathcal I_r/\mathcal I_{r-1}\arrow[r]\arrow[d,"\wr"]&0\\
0\arrow[r]&\omega_{r-1}\arrow[r]&\omega_r\arrow[r]&\omega_r/\omega_{r-1}\arrow[r]&0.
\end{tikzcd}
\]
因 $k$ 为域，两行均\Term{splitting}{分裂}：可将 $\mathcal I_{r-1}$ 的\Term{pseudobasis}{伪基} $\mathscr B_{r-1}$ 补成 $\mathcal I_r$ 的\Term{pseudobasis}{伪基} $\mathscr B_{r-1}\cup\mathscr B'_r$。于是，以 $\mathscr B'_r$ 为\Term{pseudobasis}{伪基}的\Term{closed-subspace}{闭子空间} $S_r$ 是 $\mathcal I_{r-1}$ 在 $\mathcal I_r$ 中的\Term{complementary-subspace}{补空间}，\Term{projection}{投影} $\pi:I\to\omega$ 诱导 $S_r$ 到 $\omega_{r-1}$ 在 $\omega_r$ 中的一个\Term{complementary-subspace}{补空间} $\omega'_r$ 的\Term{isomorphism}{同构}。记 $\mathcal I_\infty=\overline{\bigcup_r\mathcal I_r}$ 为\Term{closed-ideal}{闭理想}，它在 $I$ 中同样有一个\Term{complementary-subspace}{补空间} $S_\infty$；$\pi$ 诱导 $\mathcal I_\infty/\overline{I^2}$ 到 $\bigcup_r\omega_r$ 的\Term{closure}{闭包} $\omega_\infty$ 的\Term{isomorphism}{同构}，并诱导 $S_\infty$ 到 $\omega_\infty$ 在 $\omega$ 中的一个\Term{complementary-subspace}{补空间} $\omega'_\infty$ 的\Term{isomorphism}{同构}。以 $\eta$ 表示\Term{isomorphism}{同构} $S_\infty\xrightarrow{\sim}\omega'_\infty$。于是得到\Term{continuous}{连续}\Term{linear}{线性}映射
\[
\begin{tikzcd}
\overline{I^2}\times S_\infty\times\widehat{\bigoplus}_r S_r\arrow[r,"\phi"]\arrow[d,"\eta\times\theta"']&I\\
\omega=\omega'_\infty\times\omega_\infty&
\end{tikzcd}
\]
\SourcePage{597}
其中 $\widehat{\bigoplus}_rS_r$ 是 $\mathbf{PC}(k)$ 中 $S_r$ 的\Term{direct-sum}{直和}，即 $(\prod_rS_r^\dagger)^*$（见 0.2.2 的编者注（16））；$\theta:\widehat{\bigoplus}_rS_r\to\omega_\infty$ 由映射 $S_r\xrightarrow{\sim}\omega'_r\hookrightarrow\omega_\infty$ 诱导。因此，要得到所需\Term{section}{截面} $\sigma:\omega\to I$，一个充分条件（并非必要条件，见下文）是 $\theta$ 为\Term{isomorphism}{同构}。由\Term{dual}{对偶}性（见 0.2.2），这\Term{equivalence}{等价}于\Term{linear}{线性}映射 $\omega_\infty^\dagger\to\prod_rS_r^\dagger$ 为\Term{bijection}{双射}。

\textbf{5.2.1.B。} 仍记 $\omega_\infty=\overline{\bigcup_{n\in\mathbf N}\omega_n}$。所需\Term{section}{截面} $\sigma:\omega\to I$ 存在的第二种情形是：$\omega_\infty$ 有一组\Term{pseudobasis}{伪基} $\mathscr B_\infty$，它是各个 $\omega_n/\omega_{n-1}$（$n\in\mathbf N^*$）的\Term{pseudobasis}{伪基}之并（这时可再加上 $\omega/\omega_\infty$ 的一组\Term{pseudobasis}{伪基} $\mathscr B'_\infty$，得到与\Term{filtration}{滤过}相容的 $\omega$ 的\Term{pseudobasis}{伪基}）。置 $V=\omega_\infty^\dagger$，记 $V_n$ 为 $\omega_n$ 在 $V$ 中的\Term{orthogonality}{正交}\Term{space}{空间}。这\Term{equivalence}{等价}于说，在“普通”的 $k$-\Term{vector-space}{向量空间}\Term{category}{范畴}中，递降\Term{hausdorff}{分离}\Term{filtration}{滤过}
\[V=V_0\supset V_1\supset V_2\supset\cdots\]
\Term{splitting}{分裂}，即 $V$ 为满足 $F_n\simeq V_n/V_{n+1}$ 的子\Term{space}{空间} $F_n$（$n\in\mathbf N$）的\Term{direct-sum}{直和}。这未必成立：例如，若 $V$ 是元素取自 $k$ 的序列\Term{space}{空间} $\mathscr S=k^{\mathbf N}$，$\mathscr S_n$ 是满足 $i<n$ 时 $u_i=0$ 的序列 $(u_i)$ 组成的子\Term{space}{空间}，则 $\dim\mathscr S_n/\mathscr S_{n+1}=1$；但 $\mathscr S$ 并不\Term{isomorphism}{同构}于各 $\mathscr S_n/\mathscr S_{n+1}$ 的\Term{direct-sum}{直和}，因为 $\mathscr S$ 的\Term{dimension}{维数}不可数（反之，$\mathscr S$ 是这些 $\mathscr S_n/\mathscr S_{n+1}$ 的乘积，见 5.2.1.A）。然而，若 $V$ 为可数维，则上述\Term{splitting}{分裂}确实成立。\footnote{编者注（156）：本段源于与 J.-M.~Fontaine 及 E.~Bouscaren 的讨论；特别是，以下证明由 Bouscaren 告知。}事实上，取 $V$ 的一组基 $(e_n)_{n\in\mathbf N}$。对 $n$ 归纳构造递增函数 $g:\mathbf N\to\mathbf N$ 和子\Term{space}{空间} $F_i$（$i=0,\ldots,g(n)$），使 $F_i\simeq V_i/V_{i+1}$，且 $F_{\leq g(n)}=\bigoplus_{i=0}^{g(n)}F_i$ 是 $V_{g(n)+1}$ 的一个包含 $e_0,\ldots,e_n$ 的\Term{complementary-subspace}{补空间}；于是将有 $V=\bigoplus_{i\geq0}F_i$。设 $n+1\in\mathbf N$，可假定断言对 $n$ 成立（对 $n=-1$，断言为空）。若 $e_{n+1}\in F_{\leq g(n)}$，置 $g(n+1)=g(n)$。否则写成 $e_{n+1}=f+x$，其中 $f\in F_{\leq g(n)}$，$0\ne x\in V_{g(n)+1}$。令 $j$ 为使 $x\in V_j-V_{j+1}$ 的最小整数；对 $i=g(n)+1,\ldots,j$，选择 $V_{i+1}$ 在 $V_i$ 中的\Term{complementary-subspace}{补空间} $F_i$，使 $e_{n+1}\in F_j$，然后置 $g(n+1)=j$。

\textbf{5.2.1.C。}\footnote{编者注（157）：这里回到原文，并结合此前的增补作了删简。}特别地，当 $\omega$ 的\Term{filtration}{滤过}最终\Term{stable}{稳定}，即存在整数 $n_0$，使 $n_0\leq n<+\infty$ 时 $\omega_n=\omega_{n_0}$，则前述两个条件（5.2.1.A 和 B）均成立。这时 $k[[\omega]]/((x^{p^r}))_{x\in\omega_r}$ \Term{isomorphism}{同构}于下列\Term{completed-tensor-product}{完备张量积}：
\[
\frac{k[[\omega_1]]}{((x^p))_{x\in\omega_1}}
\widehat\otimes\frac{k[[\omega'_2]]}{((x^{p^2}))_{x\in\omega'_2}}
\widehat\otimes\cdots\widehat\otimes
\frac{k[[\omega'_{n_0}]]}{((x^{p^{n_0}}))_{x\in\omega'_{n_0}}}
\widehat\otimes k[[\omega'_\infty]],
\]
其中 $\omega'_n$ 是 $\omega_{n-1}$ 在 $\omega_n$ 中的\Term{complementary-subspace}{补空间}，$\omega'_\infty$ 是 $\omega_\infty=\omega_{n_0}$ 在 $\omega$ 中的\Term{complementary-subspace}{补空间}。在情形 (i)，即 $r$ \Term{large-relative-to-cardinal}{足够大}时 $\omega_r=\omega$，及情形 (ii)，即 $\omega$ \Term{finite-dimensional}{有限维}时，$\omega$ 的\Term{filtration}{滤过}显然最终\Term{stable}{稳定}；在情形 (iii)，即对每个 $r$，$I_r=0$ 时亦然（这时 $B$ 将\Term{isomorphism}{同构}于\Term{formal-power-series-algebra}{形式幂级数代数} $k[[\omega]]$）。
\SourcePage{598}
因此，结合以下 5.2.2--5.2.5，上述论述便证明了我们的\Term{theorem}{定理}。

\textbf{5.2.2。} 先设 $B$ 的\Term{height}{高度} $\leq1$，即对 $x\in I$，有 $x^p=0$。由 5.1.4，$B$ 关于\Term{filtration}{滤过} $I\supset\overline{I^2}\supset\overline{I^3}\supset\cdots$ 的\Term{associated-graded}{关联分次} $\operatorname{Gr}_I(B)$ 带有 $\mathbf{Alpg}/k$ 中的\Term{cogroup-object}{余群}结构，即\Term{profinite-algebra}{预有限}\Term{algebra}{代数} $C=\prod_{n\geq0}\operatorname{Gr}_I(B)_n$ 是一个 $k$-\Term{formal-group}{形式群} $N$ 的\Term{affine-algebra}{仿射代数}。显然 $\omega_{N/k}=I/\overline{I^2}$，且 $N$ 为\Term{height}{高度} $\leq1$ 的\Term{infinitesimal}{无穷小}\Term{formal-group}{形式群}。由 4.4，$\omega_{N/k}$ 的\Term{identity-map}{恒等映射}因而诱导一个从 $B'=k[[\omega_{N/k}]]/((x^p))_{x\in\omega_{N/k}}$ 到 $C$ 的\Term{isomorphism}{同构}。特别地，这表明若以\Term{augmentation-ideal}{增广理想}的幂\Term{filtration}{滤过} $B'$、$B$，则 5.2.1 的映射 $\phi$ 在\Term{associated-graded}{关联分次}上诱导\Term{isomorphism}{同构}。因此由 [CA]，第 V 节 7，\Term{lemma}{引理} 1（另见 5.1.5），$\phi$ 为\Term{isomorphism}{同构}。

\textbf{5.2.3。} 现设 $B$ 的\Term{height}{高度}\Term{finite}{有限}，且 $\leq r$。令 $\pi$ 为 $k$ 的\Term{isomorphism}{同构} $x\mapsto x^p$。将 $b\widehat\otimes_\pi x$ 映为 $b^px=(bx^{1/p})^p$ 的\Term{linear}{线性}映射 $B\widehat\otimes_\pi k\to B$ 有\Term{scheme-theoretic-image}{闭像}，此像正是 $B$ 的闭\Term{subalgebra}{子代数} $B^p=\{b^p\mid b\in B\}$。置 $J=B^p\cap I=B^p\cap I_A$。

记 $G_1$ 为 Frobenius \Term{morphism}{态射} $\operatorname{Fr}:G\to G^{(p)}$ 的核，$HG_1$ 为 $G^{(p)}$ 的\Term{formal-subgroup}{形式子群} $H^{(p)}$ 在 $G$ 中的\Term{inverse-image}{逆像}\Term{formal-subgroup}{形式子群}。\footnote{编者注（158）：以下内容为编者所加。}则 $HG_1$ 由 $\overline{AI}$ 中元素的 $p$ 次幂所生成的\Term{closed-ideal}{闭理想}定义，此\Term{ideal}{理想}等于 $\overline{AJ}$。另一方面，因 $G/H$ 的构造与\Term{base-change}{基变换}可交换（$G/H$ 表示\Term{flat-topology}{平坦拓扑}下的\Term{quotient-sheaf}{商层}，见 2.4），有 $(G/H)^{(p)}=G^{(p)}/H^{(p)}$，从而有\Term{commutative-diagram}{交换图}
\[
\begin{tikzcd}
G\arrow[r,"\operatorname{Fr}"]\arrow[d]&G^{(p)}\arrow[d]\\
G/H\arrow[r,"\operatorname{Fr}"']&G^{(p)}/H^{(p)}
\end{tikzcd}
\qquad
\begin{tikzcd}
A&A\widehat\otimes_\pi k\arrow[l,"{a\widehat\otimes1\mapsto a^p}"']\\
B\arrow[u]&B\widehat\otimes_\pi k\arrow[l,"{b\widehat\otimes1\mapsto b^p}"]\arrow[u]
\end{tikzcd}
\]
由此\Term{pushout}{推出}，$B^p$ 是商 $G/HG_1$ 的\Term{affine-algebra}{仿射代数}。\footnote{编者注（159）：正如原文指出的，这也可由命题 5.1 得到；但这里采用上述论证，它不使用该命题的 (i)$\Rightarrow$(ii)。}暂记 $C$ 为商 $HG_1/H$ 的\Term{affine-algebra}{仿射代数}。因 $G/H$ 的构造与\Term{base-change}{基变换}可交换，有\Term{cartesian-square}{笛卡尔方块}
\[
\begin{tikzcd}HG_1\arrow[r]\arrow[d]&G\arrow[d]\\HG_1/H\arrow[r]&G/H\end{tikzcd}
\]
从而有\Term{isomorphism}{同构} $A\widehat\otimes_B C\simeq A/\overline{AJ}=A\widehat\otimes_B(B/\overline{BJ})$。由 2.4，因 $A$、$B$ 均为\Term{local-ring}{局部环}，$A$ 在 $B$ 上\Term{topologically-free}{拓扑自由}；
\SourcePage{599}
所以自然\Term{morphism}{态射} $B/\overline{BJ}\to C$ 是\Term{isomorphism}{同构}，即 $B/\overline{BJ}$ 是 $HG_1/H$ 的\Term{affine-algebra}{仿射代数}。\footnote{编者注（159，续）：此结论亦可由命题 5.1 得到，参见上一条编者注。}当然，$B/\overline{BJ}=B_1$ 的\Term{height}{高度} $\leq1$，因为 $J=((x^p))_{x\in\mathfrak r(B)}$。

设 $B'=k[[\omega]]/((x^{p^r}))_{x\in\omega_r}$，$\phi:B'\to B$ 为 5.2.1 引入的\Term{morphism}{态射}，$B'^p$ 为\Term{subalgebra}{子代数} $\{x^p\mid x\in B'\}$，$J'$ 为 $B'^p$ 的\Term{augmentation-ideal}{增广理想}。于是有\Term{commutative-diagram}{交换图}
\[
\begin{tikzcd}
B'\arrow[r,"\phi"]\arrow[d]&B\arrow[d]\\
B'_1=B'/\overline{B'J'}\arrow[r,"\phi_1"']&B_1=B/\overline{BJ},
\end{tikzcd}
\]
且由 5.2.2，$\phi_1$ 为\Term{isomorphism}{同构}。

另一方面，由 2.4，$A$ 在 $B=\mathcal A(G/H)$ 及 $B^p=\mathcal A(G/HG_1)$ 上\Term{topologically-flat}{拓扑平坦}，故由 1.3.3，$B$ 在 $B^p$ 上\Term{topologically-flat}{拓扑平坦}。又由下面的 5.2.4，$\phi$ 诱导的\Term{morphism}{态射} $B'^p\to B^p$ 是\Term{isomorphism}{同构}。于是可以把 0.3.4 用于\Term{pseudocompact-ring}{伪紧环} $B'^p=B^p$ 及\Term{pseudocompact}{伪紧} $B^p$-模 $M=B'$、$N=B$：根据以上论证，$\phi_1=\phi\widehat\otimes_{B^p}k$ 为\Term{isomorphism}{同构}，故 $\phi$ 为\Term{isomorphism}{同构}。结合下面的 5.2.4，这就证明了 $B$ 为\Term{finite}{有限}\Term{height}{高度}时的 5.2。

\textbf{5.2.4。} 对每个\Term{pseudocompact}{伪紧} $k$-\Term{vector-space}{向量空间} $V$，以 ${}_\pi V$ 表示经标量域\Term{group-extension}{扩张} $x\mapsto x^p$ 得到的\Term{space}{空间} $V\widehat\otimes_\pi k$。\footnote{编者注（160）：因 $k$ \Term{perfect}{完美}，可把 ${}_\pi V$ 等同于\Term{abelian-group}{阿贝尔群} $V$，其中 $k$ 按 $\lambda\cdot v=\lambda^{1/p}v$ \Term{action}{作用}。}有行\Term{exact}{正合}的\Term{commutative-diagram}{交换图}
\[
\begin{tikzcd}
0\arrow[r]&{}_\pi\overline{I^2}\arrow[r,"\alpha"]\arrow[d,"u"']&{}_\pi I\arrow[r,"\beta"]\arrow[d,"v"]&{}_\pi\omega\arrow[r]\arrow[d,"w"]&0\\
0\arrow[r]&\overline{J^2}\arrow[r,"\gamma"']&J\arrow[r,"\delta"']&\overline\omega\arrow[r]&0,
\end{tikzcd}
\]
其中置 $\overline\omega=J/\overline{J^2}$，映射 $u,v,w$ 由\Term{linear}{线性}映射 ${}_\pi B\to B^p$，$x\widehat\otimes a\mapsto x^pa$，诱导。因 $u,v$ 为\Term{surjectivity}{满射}，$w$ 为\Term{surjectivity}{满射}，且其核为 ${}_\pi I_1=\operatorname{Ker}(v)$ 的像 ${}_\pi\omega_1$。

置 $J_n=\{x\in J\mid x^{p^n}=0\}$ 及 $\overline\omega_n=\delta(I_n)$，则对所有 $n\geq0$，$J_n=v({}_\pi I_{n+1})$，$\overline\omega_n=w({}_\pi\omega_{n+1})$。于是，由 5.2.1 的\Term{section}{截面} $\sigma$ 诱导的\Term{section}{截面} ${}_\pi\sigma:{}_\pi\omega\to{}_\pi I$ 通过取商定义一个\Term{section}{截面} $\tau:\overline\omega\to J$，它与 $J$、$\overline\omega$ 的\Term{filtration}{滤过}相容。因 $B^p$ 的\Term{height}{高度} $\leq r-1$，根据归纳假设，此\Term{section}{截面}诱导\Term{isomorphism}{同构}
\[
\psi:B''=k[[\overline\omega]]/((x^{p^n}))_{x\in\overline\omega_n}\xrightarrow{\sim}B^p.
\]
而 $B''$ 等同于 $B'^p$，$\psi$ 等同于 $\phi$ 诱导的\Term{morphism}{态射} $B'^p\to B^p$，故
\SourcePage{600}
当 $B$ 为\Term{finite}{有限}\Term{height}{高度}时，我们的\Term{theorem}{定理}得证。

\textbf{注记 5.2.4.A。}\footnote{编者注（161）：这条注记为编者所加，将用于 5.2.5。}设 $B$ 的\Term{height}{高度} $\leq r+1$（$r\in\mathbf N^*$），则 $I_{r+1}=I$，且有\Term{isomorphism}{同构}
\[
B\simeq
\frac{k[[S_1]]}{((x_1^p))_{x_1\in S_1}}\widehat\otimes_k\cdots\widehat\otimes_k
\frac{k[[S_r]]}{((x_r^{p^r}))_{x_r\in S_r}}\widehat\otimes_k
\frac{k[[S_{r+1}]]}{((x_{r+1}^{p^{r+1}}))_{x_{r+1}\in S_{r+1}}}.\tag{1}
\]
其中每个 $S_n$ 都是 $\overline{I^2}+I_{n-1}$ 在 $\overline{I^2}+I_n$ 中的一个\Term{complementary-subspace}{补空间}。于是 $\omega=I/\overline{I^2}$ 等同于 $\prod_{i=1}^{r+1}S_i$；容易看出，对 $n=1,\ldots,r+1$，$I_n$ 在 $\omega$ 中的像 $\omega_n$ 等同于 $\prod_{i=1}^nS_i$。

这有如下推论。设 $B_r=B/J_r$，其中 $J_r=((x^{p^r}))_{x\in\mathfrak r(B)}$，并设 $\mathfrak m=I/J_r$ 为 $B_r$ 的\Term{augmentation-ideal}{增广理想}。因 $J_r\subset\overline{I^2}$，$\omega(r)=\mathfrak m/\overline{\mathfrak m^2}$ 等同于 $\omega$。对 $n=1,\ldots,r$，记 $\omega(r)_n$ 为 $\mathfrak m_n$ 在 $\omega(r)$ 中的像；它也是 $\{x\in I\mid x^{p^n}\in J_r\}$ 在 $\omega$ 中的像，所以 $\omega(r)_n$ 包含 $\omega_n$。另一方面，\Term{isomorphism}{同构} (1) 表明 $J_r=((x_{r+1}^{p^r}))_{x_{r+1}\in S_{r+1}}$，故
\[
B_r\simeq
\frac{k[[S_1]]}{((x_1^p))_{x_1\in S_1}}\widehat\otimes_k\cdots\widehat\otimes_k
\frac{k[[S_r]]}{((x_r^{p^r}))_{x_r\in S_r}}\widehat\otimes_k
\frac{k[[S_{r+1}]]}{((x_{r+1}^{p^r}))_{x_{r+1}\in S_{r+1}}}.\tag{2}
\]
由前述论证，对 $n=1,\ldots,r-1$，$\omega(r)_n$ 等同于 $\prod_{i=1}^nS_i$。所以，对 $n=1,\ldots,r-1$，包含关系 $\omega_n\subset\omega(r)_n$ 为等号。

\textbf{5.2.5。} 还需考虑 $B$ \Term{height}{高度}无限，且\Term{projection}{投影} $I\to\omega$ 有一个与 $\omega$ 和 $I$ 的\Term{filtration}{滤过}相容的\Term{section}{截面} $\sigma$ 的情形。考虑 $\sigma$ 诱导的\Term{morphism}{态射}
\[\phi:B'=k[[\omega]]/((x^{p^n}))_{x\in\omega_n}\longrightarrow B;\]
只需证明：对所有 $r\in\mathbf N$，$\phi$ 诱导的映射 $\phi_r:B'_r\to B_r$ \Term{invertible}{可逆}。\footnote{编者注（162）：事实上，$B'$（相应地，$B$）是各 $B'_r$（相应地，$B_r$）的\Term{projective-limit}{投射极限}。此外，以下结合 5.2.3 中的增补修改了原文。}

对 $r\in\mathbf N^*$，记 $G_r$ 为迭代 Frobenius \Term{morphism}{态射} $G\to G^{(p^r)}$ 的核，$HG_r$ 为 $G^{(p^r)}$ 的\Term{formal-subgroup}{形式子群} $H^{(p^r)}$ 在 $G$ 中的\Term{inverse-image}{逆像}\Term{formal-subgroup}{形式子群}。于是 $HG_r$ 由 $\overline{AI}$ 中元素的 $p^r$ 次幂生成的\Term{closed-ideal}{闭理想}定义，该\Term{ideal}{理想}等于 $\overline{AJ_r}$，其中 $J_r=\{x^{p^r}\mid x\in I\}$。完全如 5.2.3，可知 $B_r=B/\overline{BJ_r}$ 是 $HG_r/H$ 的\Term{affine-algebra}{仿射代数}（当然其\Term{height}{高度} $\leq r$）。

记 $\mathfrak m(r)=I/\overline{BJ_r}$ 为 $B_r$ 的\Term{augmentation-ideal}{增广理想}。$B\to B_r$ 的\Term{canonical-projection}{典范投影}显然诱导 $\omega=I/\overline{I^2}$ 到 $\omega(r)=\mathfrak m(r)/\overline{\mathfrak m(r)^2}$ 的\Term{isomorphism}{同构}，从而可将二者等同。设 $\omega(r)_n$ 为\Term{closed-ideal}{闭理想} $\mathfrak m(r)_n=\{y\in\mathfrak m(r)\mid y^{p^n}=0\}$ 在 $\omega(r)$ 中的像；它也是\Term{closed-ideal}{闭理想}
\SourcePage{601}
$I(r)_n=\{x\in I\mid x^{p^n}\in\overline{BJ_r}\}$ 在 $\omega$ 中的像。\footnote{编者注（163）：以下详述了原文。}显然，当 $n\geq r$ 时，$\omega_n(r)=\omega$；下面证明当 $n<r$ 时，$\omega_n(r)=\omega_n$。对每个 $r,n$，下列序列\Term{exact}{正合}：
\[0\longrightarrow I(r)_n\cap\overline{I^2}\longrightarrow I(r)_n\longrightarrow\omega(r)_n\longrightarrow0.\]
又对固定的 $n$，$\bigcap_rI(r)_n=I_n$，因为 $\bigcap_r\overline{BJ_r}=0$。因 $\mathbf{PC}(k)$ 中\Term{filtration}{滤过}\Term{projective-limit}{投射极限}\Term{exact}{正合}（见 0.2），对每个 $n$ 有
\[\omega_n=\bigcap_r\omega(r)_n.\tag{*}\]
另一方面，由注记 5.2.4.A，当 $n<r$ 时，$\omega(r)_n=\omega(r+1)_n$。结合 (*)，得到 $n<r$ 时 $\omega(r)_n=\omega_n$。

因此，\Term{vector-space}{向量空间} $\omega$ 由子\Term{space}{空间} $(\omega(r)_n)_{n\geq0}$ 定义的\Term{filtration}{滤过}，正是 ${}_r\omega$（5.2 的记号）。从而，由 $\sigma:\omega\to I$ 与\Term{projection}{投影} $I\to\mathfrak m(r)$ 复合得到的映射 $\sigma(r)$，与 $\omega$、$\mathfrak m(r)$ 的\Term{filtration}{滤过} $(\omega(r)_n)$、$(\mathfrak m(r)_n)$ 相容。因 $k[[{}_r\omega]]/((x^{p^n}))_{x\in{}_r\omega_n}$ 正是 $B'_r$，且 $\phi_r:B'_r\to B_r$ 是 $\sigma(r)$ 诱导的\Term{morphism}{态射}，已经证明的\Term{finite}{有限}\Term{height}{高度}\Term{algebra}{代数}的结论表明 $\phi_r$ 是\Term{isomorphism}{同构}。

\begin{SGAresult}{定义}{5.2.6}
设 $(A_\lambda)_{\lambda\in\Lambda}$ 为一族\Term{profinite-algebra}{预有限} $k$-\Term{algebra}{代数}，每个都带有一个\Term{augmentation}{增广} $\varepsilon_\lambda:A_\lambda\to k$（特别地，若每个 $A_\lambda$ 都是\Term{residue-field}{剩余域}为 $k$ 的\Term{local-ring}{局部环}，则属于这一情形）。\footnote{编者注（164）：此段为编者所加，以定义 5.3 (a) 中使用的无限\Term{tensor-product}{张量积}。}定义无限\Term{completed-tensor-product}{完备张量积}
\[\mathscr A=\widehat{\bigotimes}_{\lambda\in\Lambda}A_\lambda\]
为 $\mathbf{Alp}/k$ 中 $A_F=\widehat{\bigotimes}_{\lambda\in F}A_\lambda$ 的\Term{projective-limit}{投射极限}，其中 $F$ 遍历 $\Lambda$ 的\Term{finite}{有限}\Term{subset}{子集}；当 $F'=F\cup\{\lambda\}$ 时，过渡\Term{morphism}{态射} $A_{F'}\to A_F$ 为 $\mathrm{id}\widehat\otimes\varepsilon_\lambda$。特别地，若 $\Lambda=\mathbf N^*$，并记 $X_n=\operatorname{Spf}(A_n)$ 为相应的 $k$-\Term{formal-variety}{形式簇}，则 $\operatorname{Spf}(\mathscr A)$ 表示以下\Term{functor}{函子}：它把 $C\in\mathbf{Alf}/k$ 映为 $\prod_{n\geq1}X_n(C)$ 中的“\Term{finite}{有限}”序列的\Term{set}{集合}，即
\[(x_1,x_2,\ldots)\in\prod_{n\geq1}X_n(C),\]
其中 $n$ \Term{large-relative-to-cardinal}{足够大}时 $x_n=\varepsilon_n$；这里滥用记号，以 $\varepsilon_n$ 表示 $\varepsilon_n:A_n\to k$ 与\Term{structure-morphism}{结构态射} $k\to C$ 的复合。（若每个 $A_n$ 又都是某个\Term{algebra}{代数} $k[[\omega'_n]]$ 的商，可把在 $\omega'_n$ 上为零的唯一\Term{morphism}{态射} $A_n\to C$ 记为 $0$，于是得到 $n$ \Term{large-relative-to-cardinal}{足够大}时“$x_n=0$”的序列所成的\Term{set}{集合}。）
\end{SGAresult}

\textbf{5.3。注记。} (a) 称一个\Term{profinite-algebra}{预有限} $k$-\Term{algebra}{代数}为\Term{stationary-algebra}{停驻}的（stationnaire），若它是一族\Term{truncated-formal-power-series-algebra}{截断形式幂级数代数}的\Term{completed-tensor-product}{完备张量积}。\footnote{编者注（165）：作者大概考虑的是\Term{tensor-product}{张量积} $A=\widehat{\bigotimes}_{n\in\mathbf N^*}k[[\omega'_n]]/((x^{p^n}))_{x\in\omega'_n}$，其中 $\omega'_n$ 为任意\Term{pseudocompact}{伪紧} $k$-\Term{vector-space}{向量空间}。这时不难看出，$\omega=\omega_{A/k}$ 等同于各 $\omega'_n$ 的乘积，且\Term{filtration}{滤过} $(\omega_n)_{n\in\mathbf N^*}$ 由 $\omega_n=\prod_{i=1}^n\omega'_i$ 给出，即属于 5.2.1.B 的情形。因此，把这些\Term{algebra}{代数}称为\Term{stable-algebra}{稳定}的（stables），而非“\Term{stationary-algebra}{停驻}的”，会更恰当，见 [Di73]，II，第 2.9 节，第 75 页。例如，若对所有 $n\in\mathbf N^*$，$\dim\omega'_n=1$，则 $A$ 表示把 $C\in\mathbf{Alf}/k$ 映为满足 $x_n^{p^n}=0$ 且在 $n$ \Term{large-relative-to-cardinal}{足够大}时 $x_n=0$ 的 $C$ 中元素序列 $(x_n)_{n\in\mathbf N^*}$ 所成\Term{set}{集合}的\Term{functor}{函子}。最后指出，这一情形（即 $\omega=\prod_{n\in\mathbf N^*}\omega'_n$ 的情形）对应于 M.~E.~Sweedler 在\Term{connected}{连通}\Term{cocommutative}{余交换} \Term{hopf-algebra}{Hopf 代数}的\Term{dual}{对偶}情形中所研究的情形，见 [Sw67]，\Term{theorem}{定理} 3。}
\SourcePage{602}
若 $G$ 为\Term{infinitesimal}{无穷小} $k$-\Term{formal-group}{形式群}，$B$ 为 $G$ 的一个\Term{homogeneous-space}{齐性空间}的\Term{affine-algebra}{仿射代数}，则由\Term{theorem}{定理} 5.2，对每个整数 $r\in\mathbf N$，\Term{algebra}{代数} $B/((x^{p^r}))_{x\in\mathfrak r(B)}$ 都是\Term{stationary-algebra}{停驻}的。特别地，这意味着 $B$ 是\Term{stationary-algebra}{停驻}\Term{algebra}{代数}的一个\Term{projective-limit}{投射极限}。\footnote{编者注（166）：但这样的\Term{projective-limit}{投射极限}未必是\Term{stable-algebra}{稳定}的\Term{profinite-algebra}{预有限} $k$-\Term{algebra}{代数}（按上一条编者注的意义）。例如，设 $\mathscr S$ 为由元素取自 $k$ 的序列 $(u_1,u_2,\ldots)$ 组成的“普通”$k$-\Term{vector-space}{向量空间}，$\omega=\mathscr S^*$。则 $\omega$ 是 $\mathbf{PC}(k)$ 中各个 $k$ 的副本 $k_n$（$n\in\mathbf N^*$）的\Term{direct-sum}{直和}，即属于 5.2.1.A 的情形。若 $x_n\in\omega$ 由 $x_n(\mathbf u)=u_n$（$\mathbf u=(u_i)_{i\in\mathbf N^*}$ 为任意序列）定义，则 $k$-\Term{algebra}{代数} $A=k[[\omega]]/((x_n^{p^n}))_{n\in\mathbf N^*}$ 满足 $\omega_{A/k}=\omega$、$\omega_n=\prod_{i=1}^nkx_i$，但并不\Term{stable-algebra}{稳定}：$\operatorname{Spf}(A)$ 表示把 $C\in\mathbf{Alf}/k$ 映为满足对每个 $n\in\mathbf N^*$ 都有 $x_n^{p^n}=0$ 的 $C$ 中元素的“无限”序列 $(x_n)_{n\in\mathbf N^*}$ 所成\Term{set}{集合}的\Term{functor}{函子}。}

(b) 沿用 5.2.1 的记号，我不知道是否能够选择 $A$、$B$，使得不存在与\Term{filtration}{滤过}相容的\Term{section}{截面} $\sigma:\omega\to I$。\footnote{编者注（167）：除 5.2.1.A 和 B 考虑的情形以外，编者也不知道答案。}不过，$\omega$ 可以是任意一个由递增\Term{closed-subspace}{闭子空间}列\Term{filtration}{滤过}的\Term{pseudocompact}{伪紧}\Term{vector-space}{向量空间}。事实上，若 $\omega_1\subset\omega_2\subset\cdots\subset\omega$ 为这样的\Term{filtration}{滤过}\Term{space}{空间}，则可在\Term{algebra}{代数} $B=A=k[[\omega]]/((x^{p^r}))_{x\in\omega_r}$ 中定义满足 2.1 的条件 (i)、(ii)、(iii) 的\Term{diagonal-morphism}{对角态射} $\Delta_A:A\to A\widehat\otimes_kA$；只需当 $y$ 为 $\omega$ 中元素在 $A$ 中的像时，置 $\Delta_A(y)=y\widehat\otimes1+1\widehat\otimes y$。

\begin{SGAresult}{推论}{5.4}
设 $G$ 为\Term{characteristic}{特征} $p>0$ 的\Term{perfect-field}{完美域} $k$ 上的\Term{algebraic-group}{代数群}，$H$ 为 $G$ 的\Term{algebraic-subgroup}{代数子群}，$e$ 为 $G$ 的\Term{identity-element}{单位元}在 $G/H$ 中的像，$A$ 为 $G/H$ 在 $e$ 处的\Term{local-ring}{局部环}。则 $\widehat A$ \Term{isomorphism}{同构}于下列\Term{form}{形式}的\Term{algebra}{代数}：
\[k[[X_1,\ldots,X_r,\ldots,X_s]]/(X_1^{p^{n_1}},\ldots,X_r^{p^{n_r}}).\]
\end{SGAresult}
\begin{proof}
考虑\Term{infinitesimal}{无穷小}\Term{formal-group}{形式群} $\widehat G=\operatorname{Spf}(\widehat{\mathcal O}_{G,e})$、$\widehat H=\operatorname{Spf}(\widehat{\mathcal O}_{H,e})$。由 1.3.4，$A=\mathcal O_{G/H,e}$ 的\Term{completion}{完备化} $\widehat A$ \Term{isomorphism}{同构}于 $\mathcal A(\widehat G/\widehat H)$，故推论由\Term{theorem}{定理} 5.2 (ii) 得到。\footnote{编者注（168）：另见 [DG70]，第 III.3 节，\Term{theorem}{定理} 6.1。}
\end{proof}

\textbf{5.5。补充。}\footnote{编者注（169）：本小节为编者所加，给出第 III 讲和第 VI$_A$ 讲提到的 5.4 的若干推论。}回顾以下定义。一方面，称\Term{noetherian}{诺特}\Term{local-ring}{局部环} $A$ 为\Term{complete-intersection}{完全交}，若其\Term{completion}{完备化} $\widehat A$ 是某个\Term{complete}{完备}\Term{regular}{正则}\Term{noetherian}{诺特}\Term{local-ring}{局部环} $B$ 模掉一个由 $B$ 中\Term{regular-sequence}{正则序列}生成的\Term{ideal}{理想} $I$ 的商（见 EGA IV$_4$，19.3.1）。另一方面，设 $\tau:Y\hookrightarrow X$ 为\Term{scheme}{概形}的\Term{closed-immersion}{闭浸入}。若 $y\in Y$，称 $\tau$ 在点 $y$ 为\Term{regular-immersion}{正则浸入}，若 $\mathcal O_{X,y}\to\mathcal O_{Y,y}$ 的核由一个\Term{regular-sequence}{正则序列}生成；若 $X$ 又\Term{locally-noetherian}{局部诺特}，且 $\tau$ 在每一点都是\Term{regular-immersion}{正则浸入}，则称 $\tau$ 为\Term{regular-immersion}{正则浸入}（见同上，命题 16.9.10 及定义 16.9.2）。

\SourcePage{603}
\begin{SGAresult}{推论}{5.5.1}
若 $G$ 为域 $k$ 上的\Term{algebraic-group}{代数群}，则\Term{local-ring}{局部环} $\mathcal O_{G,e}$ 为\Term{complete-intersection}{完全交}。
\end{SGAresult}
\begin{proof}
由 EGA IV$_4$，19.3.4，可设 $k$ \Term{algebraically-closed}{代数闭}。若 $\operatorname{char}(k)=0$，已知 $G$ \Term{smooth}{光滑}（见 3.3.1 或第 VI$_B$ 讲 1.6.1），故由 EGA IV$_4$，17.5.3 (d$''$)，$\mathcal O_{G,e}$ 为 $k$-\Term{formal-power-series-algebra}{形式幂级数代数}。若 $\operatorname{char}(k)=p>0$，把 5.4 用于 $H=\{e\}$，可得 $\mathcal O_{G,e}$ 为\Term{complete-intersection}{完全交}。
\end{proof}

\textbf{注记 5.5.2。} 设 $k$ 为域，$G$ 为\Term{smooth}{光滑} $k$-\Term{algebraic-group}{代数群}，$H$ 为 $G$ 的\Term{closed-subgroup}{闭子群}。

(a) 在第 III 讲 4.15 已见\Term{immersion}{浸入} $H\hookrightarrow G$ 为\Term{regular-immersion}{正则浸入}；也可由 5.4 \Term{pushout}{推出}如下。如该处，可设 $k$ \Term{algebraically-closed}{代数闭}，且只需证明 $\mathcal O_{G,e}\to\mathcal O_{H,e}$ 的核 $I$ 由一个\Term{regular-sequence}{正则序列}生成。置 $A=\mathcal O_{G,e}$、$\widehat G=\operatorname{Spf}(\widehat A)$、$\widehat H=\operatorname{Spf}(\widehat{\mathcal O}_{H,e})$。因 $A$ \Term{noetherian}{诺特}，有\Term{exact-sequence}{正合列}
\[
0\longrightarrow I\otimes_A\widehat A\longrightarrow\widehat A\xrightarrow{\pi}\mathcal A(\widehat H)\longrightarrow0,
\]
且 $\widehat A$ 在 $A$ 上\Term{faithfully-flat}{忠实平坦}。因此，由 EGA IV$_4$，16.9.10 (ii) 及 19.1.5 (ii)，只需证明 $\pi$ 的核 $\widehat I=I\otimes_A\widehat A$ 由 $\widehat A$ 中的一个\Term{regular-sequence}{正则序列}生成。

因 $G$ \Term{smooth}{光滑}，$\widehat A$ \Term{reduced}{约化}；故由 5.4，\Term{subalgebra}{子代数} $B=\mathcal A(\widehat G/\widehat H)$ \Term{isomorphism}{同构}于\Term{formal-power-series-algebra}{形式幂级数代数} $k[[x_1,\ldots,x_n]]$，从而 $\widehat G/\widehat H$ 的\Term{unit-section}{单位截面}在 $B$ 中由\Term{regular-sequence}{正则序列} $(x_1,\ldots,x_n)$ 定义。因 $\widehat A$ \Term{noetherian}{诺特}，$\widehat A$ 中由 $x_1,\ldots,x_n$ 生成的\Term{ideal}{理想} $J$ 是闭的，故依推论 1.4 等于 $\widehat I$。又因 $\widehat A$ 在 $B$ 上\Term{topologically-flat}{拓扑平坦}，从而\Term{flat}{平坦}（见 0.3.8），由 EGA IV$_4$，19.1.5 (ii)，$(x_1,\ldots,x_n)$ 是 $\widehat A$ 中的\Term{regular-sequence}{正则序列}。

(b) 还可由 5.2 (ii) \Term{pushout}{推出}以下更精确的结果。设 $k$ \Term{perfect}{完美}。把 5.2 (ii) 用于\Term{algebra}{代数} $C=\mathcal A(\widehat H)$，可得 $\omega_H$ 的一组基 $(y_1,\ldots,y_{r+s})$ 及整数 $1\leq n_1\leq\cdots\leq n_r$，使 $\mathcal A(\widehat H)$ \Term{isomorphism}{同构}于 $k[[y_1,\ldots,y_{r+s}]]$ 模掉由 $y_i^{p^{n_i}}$（$i=1,\ldots,r$）生成的\Term{ideal}{理想}所得的商。把 $y_i$ \Term{lifting}{提升}为 $\omega_G$ 中的元素 $x_i$，并把 $(x_1,\ldots,x_{r+s})$ 补成 $\omega_G$ 的基 $(x_1,\ldots,x_n)$。因 $\mathcal A(\widehat G)$ \Term{reduced}{约化}，由 5.2 (iii)，\Term{morphism}{态射} $k[[x_1,\ldots,x_n]]\to\mathcal A(\widehat G)$ 为\Term{isomorphism}{同构}。于是得到：$\widehat G$ 存在一个“坐标系”$(x_1,\ldots,x_n)$（即\Term{isomorphism}{同构} $\mathcal A(\widehat G)\simeq k[[x_1,\ldots,x_n]]$），使得 $\widehat H$ 由方程
\[x_i^{p^{n_i}}=0\quad(i=1,\ldots,r),\qquad x_i=0\quad(i>r+s)\]
定义（“Dieudonn\'e \Term{theorem}{定理}”，比较 [Di55]，第 19 节，\Term{theorem}{定理} 6，及 [Di73]，II，第 3.2 节，命题 3 及其前文）。

\SGAsection{}{参考文献}
\begin{description}
\item[{[CA]}] P.~Gabriel，《论\Term{abelian-category}{阿贝尔范畴}》（\textit{Des catégories abéliennes}），Bull. Soc. Math. France \textbf{90} (1962)，323--448。\footnote{编者注（170）：原文仅列此文献，其余文献为编者补充。}
\SourcePage{604}
\item[{[Ab80]}] E.~Abe，《\Term{hopf-algebra}{Hopf 代数}》（\textit{Hopf algebras}），Cambridge Univ. Press，1980。
\item[{[Br00]}] C.~Breuil，《$p$-\Term{divisible-group}{可除群}、\Term{finite-group}{有限群}与\Term{filtered}{滤过}模》（\textit{Groupes $p$-divisibles, groupes finis et modules filtrés}），Ann. of Math. \textbf{152} (2000)，第 2 期，489--549。
\item[{[BAlg]}] N.~Bourbaki，《\Term{algebra}{代数}》（\textit{Algèbre}），第 I--III 章及第 IV--VII 章，Hermann，1970；Masson，1981。
\item[{[BAC]}] N.~Bourbaki，《\Term{commutative-algebra}{交换代数}》（\textit{Algèbre commutative}），第 I--IV 章，Masson，1985。
\item[{[BEns]}] N.~Bourbaki，《\Term{set-theory}{集合论}》（\textit{Théorie des ensembles}），第 I--IV 章，Hermann，1970。
\item[{[BLie]}] N.~Bourbaki，《Lie 群与 \Term{lie-algebra}{Lie 代数}》（\textit{Groupes et algèbres de Lie}），第 II--III 章，Hermann，1970。
\item[{[Ca62]}] P.~Cartier，《\Term{algebraic-group}{代数群}与\Term{formal-group}{形式群}》（\textit{Groupes algébriques et groupes formels}），载《\Term{algebraic-group}{代数群}理论讨论会》（布鲁塞尔，1962），第 87--111 页，Gauthier-Villars，1962。
\item[{[DG70]}] M.~Demazure、P.~Gabriel，《\Term{algebraic-group}{代数群}》（\textit{Groupes algébriques}），Masson \& North-Holland，1970。
\item[{[De72]}] M.~Demazure，《$p$-\Term{divisible-group}{可除群}讲义》（\textit{Lectures on $p$-divisible groups}），Lect. Notes Math. \textbf{302}，Springer-Verlag，1972。
\item[{[Di55]}] J.~Dieudonn\'e，《\Term{characteristic}{特征} $p>0$ 的域上的 Lie 群与 Lie 超\Term{algebra}{代数}》（\textit{Groupes de Lie et hyperalgèbres de Lie sur un corps de caractéristique $p>0$}），Comm. Math. Helv. \textbf{28} (1955)，87--118。
\item[{[Di73]}] J.~Dieudonn\'e，《\Term{formal-group}{形式群}理论导论》（\textit{Introduction to the theory of formal groups}），Marcel Dekker，1973。
\item[{[Fo77]}] J.-M.~Fontaine，《\Term{local-scheme}{局部}域上的 $p$-\Term{divisible-group}{可除群}》（\textit{Groupes $p$-divisibles sur les corps locaux}），Astérisque \textbf{47--48}，Soc. Math. France，1977。
\item[{[Gr57]}] A.~Grothendieck，《\Term{homological-algebra}{同调代数}的若干问题》（\textit{Sur quelques points d'algèbre homologique}），Tôhoku Math. J. \textbf{9} (1957)，119--221。
\item[{[Gr74]}] A.~Grothendieck，《Barsotti--Tate 群与 Dieudonn\'e 晶体》（\textit{Groupes de Barsotti-Tate et cristaux de Dieudonné}），Presses Univ. Montréal，1974。
\item[{[HS69]}] R.~G.~Heyneman、M.~E.~Sweedler，《\Term{affine}{仿射} \Term{hopf-algebra}{Hopf 代数} I》（\textit{Affine Hopf Algebras I}），J. Algebra \textbf{13} (1969)，194--241。
\item[{[Il85]}] L.~Illusie，《Barsotti--Tate 群的\Term{deformation}{形变}，据 A.~Grothendieck》（\textit{Déformations de groupes de Barsotti-Tate, d'après A. Grothendieck}），载《算术\Term{pencil}{铅笔}讨论班：Mordell 猜想》，第 151--198 页，Astérisque \textbf{127}，Soc. Math. France，1985。
\item[{[Ja65]}] I.~M.~James，Math. Reviews 中对 [MM65] 的评论，MR0174052。
\item[{[La75]}] M.~Lazard，《交换\Term{formal-group}{形式群}》（\textit{Commutative formal groups}），Lect. Notes Math. \textbf{443}，Springer-Verlag，1975。
\item[{[LT65]}] J.~Lubin、J.~Tate，《\Term{local-scheme}{局部}域中的\Term{form}{形式}复乘》（\textit{Formal complex multiplication in local fields}），Ann. of Math. \textbf{81} (1965)，380--387。
\item[{[LT66]}] J.~Lubin、J.~Tate，《单参数\Term{form}{形式} Lie 群的\Term{form}{形式}模》（\textit{Formal moduli for one-parameter formal Lie groups}），Bull. Soc. Math. France \textbf{94} (1966)，49--60。
\item[{[Ma91]}] A.~Masuoka，《具有\Term{cocommutative}{余交换}\Term{coradical}{余根基}的 \Term{hopf-algebra}{Hopf 代数}》（\textit{On Hopf algebras with cocommutative coradicals}），J. Algebra \textbf{144} (1991)，451--466。
\item[{[Me72]}] W.~Messing，《与 Barsotti--Tate 群相联系的晶体及其在\Term{abelian-scheme}{阿贝尔概形}中的应用》（\textit{The crystals associated with Barsotti-Tate Groups: with applications to abelian schemes}），Lect. Notes Math. \textbf{264}，Springer-Verlag，1972。
\item[{[MM65]}] J.~W.~Milnor、J..~C.~Moore，《论 \Term{hopf-algebra}{Hopf 代数}的结构》（\textit{On the structure of Hopf algebras}），Ann. of Math. \textbf{81} (1965)，211--264。
\item[{[Mi65]}] B.~Mitchell，《\Term{category}{范畴}理论》（\textit{Theory of categories}），Academic Press，1965。
\SourcePage{605}
\item[{[Ne75]}] K.~Newman，《\Term{cocommutative}{余交换} \Term{hopf-algebra}{Hopf 代数}中的双\Term{ideal}{理想}与子 \Term{hopf-algebra}{Hopf 代数}之间的对应》（\textit{A correspondence between bi-ideals and sub-Hopf algebras in cocommutative Hopf algebras}），J. Algebra \textbf{36} (1975)，1--15。
\item[{[Po73]}] N.~Popescu，《\Term{abelian-category}{阿贝尔范畴}及其在环与模中的应用》（\textit{Abelian categories with applications to rings and modules}），Academic Press，1973。
\item[{[Sch90]}] H.~J.~Schneider，《任意 \Term{hopf-algebra}{Hopf 代数}的\Term{principal-homogeneous-space}{主齐性空间}》（\textit{Principal homogeneous spaces for arbitrary Hopf algebras}），Israel J. Math. \textbf{72} (1990)，第 1--2 期，167--195。
\item[{[Sw67]}] M.~Sweedler，《只有一个群状元素的 \Term{hopf-algebra}{Hopf 代数}》（\textit{Hopf algebras with one grouplike element}），Trans. Amer. Math. Soc. \textbf{127} (1967)，第 3 期，515--526。
\item[{[Sw69]}] M.~Sweedler，《\Term{hopf-algebra}{Hopf 代数}》（\textit{Hopf algebras}），Benjamin，1969。
\item[{[Tak72]}] M.~Takeuchi，《Hopf \Term{ideal}{理想}与子 \Term{hopf-algebra}{Hopf 代数}之间的对应》（\textit{A correspondence between Hopf ideals and sub-Hopf algebras}），Manuscripta Math. \textbf{7} (1972)，251--270。
\item[{[Tak79]}] M.~Takeuchi，《相对 Hopf 模：\Term{equivalence}{等价}与自由性判据》（\textit{Relative Hopf modules--Equivalence and freeness criteria}），J. Algebra \textbf{60} (1979)，452--471。
\item[{[Ta67]}] J.~Tate，《$p$-\Term{divisible-group}{可除群}》（\textit{$p$-divisible groups}），载 T.~A.~Springer 编《\Term{local-scheme}{局部}域会议录》（Driebergen，1966），第 158--183 页，Springer-Verlag，1967。
\item[{[We95]}] C.~A.~Weibel，《\Term{homological-algebra}{同调代数}导论》（\textit{An introduction to homological algebra}），Cambridge Univ. Press，1995。
\end{description}
\SourcePage{606}
% 原 PDF 第 606 页为空白页。
